% Môn: Vật lý
% Lớp: 12
% Chương: Chương II. Khí lí tưởng
% Bài: L12C2 Bài 12. Áp suất khí theo mô hình động học phân tử
% Số câu: 146
% App đọc file này trực tiếp — sửa rồi Commit trên GitHub, bấm Đồng bộ GitHub .tex

% L12C2 Bài 12. Áp suất khí theo mô hình động học phân tử

% ===== Câu 1 =====
% ID: L12C2B12-01-DS
% Mức: TH
\dangbt{Bài toán thực tế về áp suất khí theo mô hình động học phân tử}
\begin{ex}
% ID: L12C2B12-01-DS
% Mức: TH
Xét dạng “Bài toán thực tế về áp suất khí theo mô hình động học phân tử”: a) Công thức đặc trưng là $p=\frac13\rho v_{rms}^2=nkT$. b) Có thể bỏ qua điều kiện khí loãng, cân bằng nhiệt và áp suất tuyệt đối. c) p tăng khi mật độ phân tử hoặc nhiệt độ tăng. d) Có thể dùng áp suất dư thay cho áp suất tuyệt đối.
\choiceTF

\loigiai{
a) Đúng: phù hợp với công thức và điều kiện đã nêu. b) Sai: phát biểu không thỏa công thức hoặc điều kiện áp dụng. c) Đúng: phù hợp với công thức và điều kiện đã nêu. d) Sai: phát biểu không thỏa công thức hoặc điều kiện áp dụng.
}
\end{ex}

% ===== Câu 2 =====
% ID: L12C2B12-01-TLN
% Mức: TH
\begin{ex}
% ID: L12C2B12-01-TLN
% Mức: TH
Một khí có $\rho=1,2$ kg/m³ và $v_{rms}=500$ m/s. Tính áp suất theo kPa.
\shortans{100}
\loigiai{
$p=\rho v_{rms}^2/3=1,2\cdot500^2/3=100000Pa =100\,\text{kPa}$.
}
\end{ex}

% ===== Câu 3 =====
% ID: L12C2B12-01-TN
% Mức: NB
\begin{ex}
% ID: L12C2B12-01-TN
% Mức: NB
Công thức cốt lõi phù hợp nhất cho dạng “Bài toán thực tế về áp suất khí theo mô hình động học phân tử” là
\choice
{\True $p=\frac13\rho v_{rms}^2=nkT$}
{$Q=mc\Delta T$}
{$A=Fs$}
{$U=RI$}
\loigiai{
Đáp án A. Dùng $p=\frac13\rho v_{rms}^2=nkT$ trong điều kiện khí loãng, cân bằng nhiệt và áp suất tuyệt đối.
}
\end{ex}

% ===== Câu 4 =====
% ID: L12C2B12-02-TL
% Mức: TH
\begin{ex}
% ID: L12C2B12-02-TL
% Mức: TH
Nêu quy trình giải một bài thuộc dạng “Bài toán thực tế về áp suất khí theo mô hình động học phân tử”.
\choiceTF

\loigiai{
Bước 1: tóm tắt và đổi đơn vị. Bước 2: nhận diện khí loãng, cân bằng nhiệt và áp suất tuyệt đối. Bước 3: dùng $p=\frac13\rho v_{rms}^2=nkT$. Bước 4: giải, ghi đơn vị và đối chiếu p tăng khi mật độ phân tử hoặc nhiệt độ tăng.
}
\end{ex}

% ===== Câu 5 =====
% ID: L12C2B12-02-TLN
% Mức: TH
\begin{ex}
% ID: L12C2B12-02-TLN
% Mức: TH
Đổi 27 ° $C$ sang nhiệt độ tuyệt đối theo kelvin.
\shortans{300}
\loigiai{
$T=27+273=300$ K.
}
\end{ex}

% ===== Câu 6 =====
% ID: L12C2B12-02-TN
% Mức: NB
\begin{ex}
% ID: L12C2B12-02-TN
% Mức: NB
Điều kiện quan trọng khi áp dụng dạng “Bài toán thực tế về áp suất khí theo mô hình động học phân tử” là
\choice
{khối lượng vật luôn bằng $1\,\text{kg}$}
{\True khí loãng, cân bằng nhiệt và áp suất tuyệt đối}
{nhiệt độ luôn bằng 0 ° $C$}
{áp suất luôn bằng áp suất khí quyển}
\loigiai{
Đáp án B. Điều kiện áp dụng là khí loãng, cân bằng nhiệt và áp suất tuyệt đối.
}
\end{ex}

% ===== Câu 7 =====
% ID: L12C2B12-03-TL
% Mức: VD
\begin{ex}
% ID: L12C2B12-03-TL
% Mức: VD
Một khí có $\rho=1,2$ kg/m³ và $v_{rms}=500$ m/s. Tính áp suất theo kPa.
\choiceTF

\loigiai{
$p=\rho v_{rms}^2/3=1,2\cdot500^2/3=100000Pa =100\,\text{kPa}$.
}
\end{ex}

% ===== Câu 8 =====
% ID: L12C2B12-03-TLN
% Mức: VD
\begin{ex}
% ID: L12C2B12-03-TLN
% Mức: VD
Một đại lượng tăng từ 100 lên 150. Tính tỉ số giá trị sau so với trước.
\shortans{1,5}
\loigiai{
Tỉ số bằng $150/100=1,5$.
}
\end{ex}

% ===== Câu 9 =====
% ID: L12C2B12-03-TN
% Mức: TH
\begin{ex}
% ID: L12C2B12-03-TN
% Mức: TH
Nhận xét đúng về quan hệ các đại lượng là
\choice
{các đại lượng không phụ thuộc nhau}
{mọi đại lượng đều không đổi}
{\True p tăng khi mật độ phân tử hoặc nhiệt độ tăng}
{chỉ phụ thuộc thời gian}
\loigiai{
Đáp án C. p tăng khi mật độ phân tử hoặc nhiệt độ tăng.
}
\end{ex}

% ===== Câu 10 =====
% ID: L12C2B12-04-DS
% Mức: VD
\begin{ex}
% ID: L12C2B12-04-DS
% Mức: VD
Xét dạng “Bài toán thực tế về áp suất khí theo mô hình động học phân tử”: a) Quan hệ cần dùng là p tăng khi mật độ phân tử hoặc nhiệt độ tăng. b) Có thể suy ra quan hệ mà không xét điều kiện quá trình. c) Kết quả phải phù hợp chiều biến thiên vật lí. d) Áp suất dư luôn thay được áp suất tuyệt đối.
\choiceTF

\loigiai{
a) Đúng: phù hợp với công thức và điều kiện đã nêu. b) Sai: phát biểu không thỏa công thức hoặc điều kiện áp dụng. c) Đúng: phù hợp với công thức và điều kiện đã nêu. d) Sai: phát biểu không thỏa công thức hoặc điều kiện áp dụng.
}
\end{ex}

% ===== Câu 11 =====
% ID: L12C2B12-04-TLN
% Mức: VD
\begin{ex}
% ID: L12C2B12-04-TLN
% Mức: VD
Một thể tích 2,5 $L$ bằng bao nhiêu m³? Viết kết quả theo dạng $a\cdot10^{-3}$.
\shortans{2,5}
\loigiai{
$2,5L=2,5\cdot10^{-3}$ m³.
}
\end{ex}

% ===== Câu 12 =====
% ID: L12C2B12-04-TN
% Mức: TH
\begin{ex}
% ID: L12C2B12-04-TN
% Mức: TH
Sai lầm cần tránh khi giải dạng này là
\choice
{ghi đầy đủ dữ kiện}
{đổi đơn vị đúng}
{kiểm tra điều kiện}
{\True dùng áp suất dư thay cho áp suất tuyệt đối}
\loigiai{
Đáp án D. Cần tránh: dùng áp suất dư thay cho áp suất tuyệt đối.
}
\end{ex}

% ===== Câu 13 =====
% ID: L12C2B12-05-DS
% Mức: VDC
\begin{ex}
% ID: L12C2B12-05-DS
% Mức: VDC
Xét dạng “Bài toán thực tế về áp suất khí theo mô hình động học phân tử”: a) Dạng này có thể yêu cầu phối hợp đổi đơn vị và lập tỉ số. b) Kelvin được đổi sang Celsius bằng cách cộng 273. c) Cần ghi rõ đại lượng cần tìm. d) Công thức nào có đủ kí hiệu cũng dùng được.
\choiceTF

\loigiai{
a) Đúng: phù hợp với công thức và điều kiện đã nêu. b) Sai: phát biểu không thỏa công thức hoặc điều kiện áp dụng. c) Đúng: phù hợp với công thức và điều kiện đã nêu. d) Sai: phát biểu không thỏa công thức hoặc điều kiện áp dụng.
}
\end{ex}

% ===== Câu 14 =====
% ID: L12C2B12-05-TLN
% Mức: VD
\begin{ex}
% ID: L12C2B12-05-TLN
% Mức: VD
Áp suất $120\,\text{kPa}$ bằng bao nhiêu Pa, lấy hệ số của $10^3$ làm đáp án?
\shortans{120}
\loigiai{
$120\,\text{kPa}=120\cdot10^3$ Pa.
}
\end{ex}

% ===== Câu 15 =====
% ID: L12C2B12-05-TN
% Mức: TH
\begin{ex}
% ID: L12C2B12-05-TN
% Mức: TH
Đơn vị thường dùng sau khi chuẩn hóa theo SI là
\choice
{\True Pa}
{kg·m}
{W/s}
{C/mol}
\loigiai{
Đáp án A. Các đơn vị phù hợp: Pa.
}
\end{ex}

% ===== Câu 16 =====
% ID: L12C2B12-06-DS
% Mức: TH
\dangbt{Công thức áp suất theo khối lượng riêng và tốc độ phân tử}
\begin{ex}
% ID: L12C2B12-06-DS
% Mức: TH
Xét các phát biểu sau trong dạng “Công thức áp suất theo khối lượng riêng và tốc độ phân tử”.
\choiceTF
{\True Áp suất chất khí lên thành bình có nguồn gốc từ — phát biểu đúng là: các va chạm của phân tử khí với thành bình.}
{Áp suất chất khí lên thành bình có nguồn gốc từ — phát biểu trọng lượng riêng của bình là đúng.}
{\True Công thức áp suất theo mô hình động học phân tử là — phát biểu đúng là: $p=\frac13\rho\overline{v^2}$.}
{Công thức áp suất theo mô hình động học phân tử là — phát biểu $p=\rho/3v^2$ là đúng.}
\loigiai{
\begin{itemchoice}
\item (Đúng) Áp suất chất khí lên thành bình có nguồn gốc từ — phát biểu đúng là: các va chạm của phân tử khí với thành bình. (Vì): Phân tử đổi động lượng khi va chạm và gây lực lên thành bình. b) Sai: trái với hệ thức hoặc điều kiện đã nêu. c) Đúng: Mô hình động học cho $p=\frac13\rho\overline{v^2}$. d) Sai: trái với hệ thức hoặc điều kiện đã nêu
\item (Sai) Áp suất chất khí lên thành bình có nguồn gốc từ — phát biểu trọng lượng riêng của bình là đúng.
\item (Đúng) Công thức áp suất theo mô hình động học phân tử là — phát biểu đúng là: $p=\frac13\rho\overline{v^2}$.
\item (Sai) Công thức áp suất theo mô hình động học phân tử là — phát biểu $p=\rho/3v^2$ là đúng.
\end{itemchoice}
}
\end{ex}

% ===== Câu 17 =====
% ID: L12C2B12-06-TL
% Mức: VDC
\dangbt{Bài toán thực tế về áp suất khí theo mô hình động học phân tử}
\begin{ex}
% ID: L12C2B12-06-TL
% Mức: VDC
Đề xuất cách giải khi bài toán có nhiều giai đoạn liên quan đến dạng “Bài toán thực tế về áp suất khí theo mô hình động học phân tử”.
\choiceTF

\loigiai{
Chia quá trình thành các trạng thái liên tiếp, ghi p, $V$, $T$ cho từng trạng thái; áp dụng $p=\frac13\rho v_{rms}^2=nkT$ ở đoạn thỏa khí loãng, cân bằng nhiệt và áp suất tuyệt đối; dùng kết quả đoạn trước làm dữ kiện đoạn sau và kiểm tra tính liên tục.
}
\end{ex}

% ===== Câu 18 =====
% ID: L12C2B12-06-TLN
% Mức: VDC
\begin{ex}
% ID: L12C2B12-06-TLN
% Mức: VDC
Theo quan hệ “p tăng khi mật độ phân tử hoặc nhiệt độ tăng”, nếu đại lượng phụ thuộc tỉ lệ thuận tăng gấp 4 thì hệ số tăng là bao nhiêu?
\shortans{4}
\loigiai{
Với quan hệ tỉ lệ thuận, hệ số biến đổi được giữ nguyên: 4.
}
\end{ex}

% ===== Câu 19 =====
% ID: L12C2B12-06-TN
% Mức: VD
\begin{ex}
% ID: L12C2B12-06-TN
% Mức: VD
Bước đầu tiên hợp lí khi giải bài là
\choice
{thay số ngay}
{\True xác định quá trình và đại lượng không đổi}
{bỏ qua đơn vị}
{đổi áp suất thành nhiệt độ}
\loigiai{
Đáp án B. Trước hết phải nhận diện quá trình, liệt kê dữ kiện và đại lượng không đổi.
}
\end{ex}

% ===== Câu 20 =====
% ID: L12C2B12-07-DS
% Mức: TH
\dangbt{Công thức áp suất theo khối lượng riêng và tốc độ phân tử}
\begin{ex}
% ID: L12C2B12-07-DS
% Mức: TH
Xét các phát biểu sau trong dạng “Công thức áp suất theo khối lượng riêng và tốc độ phân tử”.
\choiceTF
{\True Khí có $\rho=1,4$ kg/m³ và tốc độ căn quân phương $480\,\text{m}$ /s. Áp suất gần bằng — phát biểu đúng là: $107,52\,\text{kPa}$.}
{Khí có $\rho=1,4$ kg/m³ và tốc độ căn quân phương $480\,\text{m}$ /s. Áp suất gần bằng — phát biểu $322,56\,\text{kPa}$ là đúng.}
{\True Khi nhiệt độ tuyệt đối tăng gấp 4, tốc độ căn quân phương của phân tử — phát biểu đúng là: tăng gấp 2.}
{Khi nhiệt độ tuyệt đối tăng gấp 4, tốc độ căn quân phương của phân tử — phát biểu giảm một nửa là đúng.}
\loigiai{
\begin{itemchoice}
\item (Đúng) Khí có $\rho=1,4$ kg/m³ và tốc độ căn quân phương $480\,\text{m}$ /s. Áp suất gần bằng — phát biểu đúng là: $107,52\,\text{kPa}$. (Vì): $p=\rho v_{rms}^2/3=107,52$ kPa. b) Sai: trái với hệ thức hoặc điều kiện đã nêu. c) Đúng: $v_{rms}\sim\sqrt{T}$. d) Sai: trái với hệ thức hoặc điều kiện đã nêu
\item (Sai) Khí có $\rho=1,4$ kg/m³ và tốc độ căn quân phương $480\,\text{m}$ /s. Áp suất gần bằng — phát biểu $322,56\,\text{kPa}$ là đúng.
\item (Đúng) Khi nhiệt độ tuyệt đối tăng gấp 4, tốc độ căn quân phương của phân tử — phát biểu đúng là: tăng gấp 2.
\item (Sai) Khi nhiệt độ tuyệt đối tăng gấp 4, tốc độ căn quân phương của phân tử — phát biểu giảm một nửa là đúng.
\end{itemchoice}
}
\end{ex}

% ===== Câu 21 =====
% ID: L12C2B12-07-TL
% Mức: TH
\begin{ex}
% ID: L12C2B12-07-TL
% Mức: TH
Trong dạng “Công thức áp suất theo khối lượng riêng và tốc độ phân tử”, Trình bày cơ sở lí thuyết của dạng “Công thức áp suất theo khối lượng riêng và tốc độ phân tử” và nêu điều kiện áp dụng.
\choiceTF

\loigiai{
Nêu đúng đại lượng không đổi, hệ thức đặc trưng và yêu cầu dùng nhiệt độ kelvin, áp suất tuyệt đối khi có liên quan.
}
\end{ex}

% ===== Câu 22 =====
% ID: L12C2B12-07-TLN
% Mức: TH
\begin{ex}
% ID: L12C2B12-07-TLN
% Mức: TH
Trong dạng “Công thức áp suất theo khối lượng riêng và tốc độ phân tử”, Khí có $\rho=1,2$ kg/m³, $v_{rms}=300$ m/s. Tính p theo kPa.
\shortans{36}
\loigiai{
$p=\rho v_{rms}^2/3=36$ kPa.
}
\end{ex}

% ===== Câu 23 =====
% ID: L12C2B12-07-TN
% Mức: VD
\dangbt{Bài toán thực tế về áp suất khí theo mô hình động học phân tử}
\begin{ex}
% ID: L12C2B12-07-TN
% Mức: VD
Với dữ kiện của bài mẫu: Một khí có $\rho=1,2$ kg/m³ và $v_{rms}=500$ m/s. Tính áp suất theo kPa.
\choice
{\True Kết quả 100}
{Kết quả 200}
{Không đủ dữ kiện}
{Kết quả bằng 0}
\loigiai{
Đáp án A. $p=\rho v_{rms}^2/3=1,2\cdot500^2/3=100000Pa =100\,\text{kPa}$.
}
\end{ex}

% ===== Câu 24 =====
% ID: L12C2B12-08-TL
% Mức: TH
\dangbt{Công thức áp suất theo khối lượng riêng và tốc độ phân tử}
\begin{ex}
% ID: L12C2B12-08-TL
% Mức: TH
Trong dạng “Công thức áp suất theo khối lượng riêng và tốc độ phân tử”, Mô tả dạng đồ thị đặc trưng và cách nhận biết quá trình trong dạng “Công thức áp suất theo khối lượng riêng và tốc độ phân tử”.
\choiceTF

\loigiai{
Xác định các trục, đại lượng không đổi và suy ra hình dạng đồ thị từ hệ thức vật lí tương ứng.
}
\end{ex}

% ===== Câu 25 =====
% ID: L12C2B12-08-TLN
% Mức: TH
\begin{ex}
% ID: L12C2B12-08-TLN
% Mức: TH
Trong dạng “Công thức áp suất theo khối lượng riêng và tốc độ phân tử”, Nhiệt độ tuyệt đối tăng 4 lần. Tốc độ căn quân phương tăng bao nhiêu lần?
\shortans{2}
\loigiai{
$v_{rms}\sim\sqrt T$, nên tăng 2 lần.
}
\end{ex}

% ===== Câu 26 =====
% ID: L12C2B12-08-TN
% Mức: VD
\dangbt{Bài toán thực tế về áp suất khí theo mô hình động học phân tử}
\begin{ex}
% ID: L12C2B12-08-TN
% Mức: VD
Khi kiểm tra kết quả, tiêu chí vật lí quan trọng là
\choice
{chỉ cần số dương}
{không cần đơn vị}
{\True phù hợp chiều biến thiên theo định luật}
{luôn phải lớn hơn giá trị đầu}
\loigiai{
Đáp án C. Kết quả phải phù hợp với nhận xét: p tăng khi mật độ phân tử hoặc nhiệt độ tăng.
}
\end{ex}

% ===== Câu 27 =====
% ID: L12C2B12-09-DS
% Mức: VD
\dangbt{Công thức áp suất theo khối lượng riêng và tốc độ phân tử}
\begin{ex}
% ID: L12C2B12-09-DS
% Mức: VD
Xét các phát biểu sau trong dạng “Công thức áp suất theo khối lượng riêng và tốc độ phân tử”.
\choiceTF
{\True Ở cùng nhiệt độ, phân tử khí có khối lượng nhỏ hơn có tốc độ căn quân phương — phát biểu đúng là: lớn hơn.}
{Ở cùng nhiệt độ, phân tử khí có khối lượng nhỏ hơn có tốc độ căn quân phương — phát biểu nhỏ hơn là đúng.}
{\True Nếu giữ mật độ phân tử không đổi và tăng $T$ gấp đôi thì áp suất — phát biểu đúng là: tăng gấp đôi.}
{Nếu giữ mật độ phân tử không đổi và tăng $T$ gấp đôi thì áp suất — phát biểu không đổi là đúng.}
\loigiai{
\begin{itemchoice}
\item (Đúng) Ở cùng nhiệt độ, phân tử khí có khối lượng nhỏ hơn có tốc độ căn quân phương — phát biểu đúng là: lớn hơn. (Vì): $v_{rms}=\sqrt{3kT/m}$. b) Sai: trái với hệ thức hoặc điều kiện đã nêu. c) Đúng: Từ $p=nkT$, p tỉ lệ thuận T. d) Sai: trái với hệ thức hoặc điều kiện đã nêu
\item (Sai) Ở cùng nhiệt độ, phân tử khí có khối lượng nhỏ hơn có tốc độ căn quân phương — phát biểu nhỏ hơn là đúng.
\item (Đúng) Nếu giữ mật độ phân tử không đổi và tăng $T$ gấp đôi thì áp suất — phát biểu đúng là: tăng gấp đôi.
\item (Sai) Nếu giữ mật độ phân tử không đổi và tăng $T$ gấp đôi thì áp suất — phát biểu không đổi là đúng.
\end{itemchoice}
}
\end{ex}

% ===== Câu 28 =====
% ID: L12C2B12-09-TL
% Mức: VD
\begin{bt}
% ID: L12C2B12-09-TL
% Mức: VD
Trong dạng “Công thức áp suất theo khối lượng riêng và tốc độ phân tử”, Mật độ phân tử tăng 3 lần, nhiệt độ không đổi. Áp suất tăng bao nhiêu lần?
\loigiai{
Từ $p=nkT$, p tăng 3 lần.
}
\end{bt}

% ===== Câu 29 =====
% ID: L12C2B12-09-TLN
% Mức: VD
\begin{ex}
% ID: L12C2B12-09-TLN
% Mức: VD
Trong dạng “Công thức áp suất theo khối lượng riêng và tốc độ phân tử”, Mật độ phân tử tăng 3 lần, nhiệt độ không đổi. Áp suất tăng bao nhiêu lần?
\shortans{3}
\loigiai{
Từ $p=nkT$, p tăng 3 lần.
}
\end{ex}

% ===== Câu 30 =====
% ID: L12C2B12-09-TN
% Mức: VDC
\dangbt{Bài toán thực tế về áp suất khí theo mô hình động học phân tử}
\begin{ex}
% ID: L12C2B12-09-TN
% Mức: VDC
Nếu giữ đúng các điều kiện của định luật, biểu thức nên dùng là
\choice
{$s=vt$}
{\True $p=\frac13\rho v_{rms}^2=nkT$}
{$P=A/t$}
{$F=ma$}
\loigiai{
Đáp án B. Biểu thức đặc trưng là $p=\frac13\rho v_{rms}^2=nkT$.
}
\end{ex}

% ===== Câu 31 =====
% ID: L12C2B12-10-DS
% Mức: VDC
\dangbt{Công thức áp suất theo khối lượng riêng và tốc độ phân tử}
\begin{ex}
% ID: L12C2B12-10-DS
% Mức: VDC
Xét các phát biểu sau trong dạng “Công thức áp suất theo khối lượng riêng và tốc độ phân tử”.
\choiceTF
{\True Khi va chạm đàn hồi với thành bình đứng yên, đại lượng của phân tử đổi dấu là — phát biểu đúng là: thành phần vận tốc vuông góc thành.}
{Khi va chạm đàn hồi với thành bình đứng yên, đại lượng của phân tử đổi dấu là — phát biểu khối lượng là đúng.}
{\True Đơn vị của hằng số Boltzmann k là — phát biểu đúng là: J/K.}
{Đơn vị của hằng số Boltzmann k là — phát biểu Pa/ $K$ là đúng.}
\loigiai{
\begin{itemchoice}
\item (Đúng) Khi va chạm đàn hồi với thành bình đứng yên, đại lượng của phân tử đổi dấu là — phát biểu đúng là: thành phần vận tốc vuông góc thành. (Vì): Thành phần pháp tuyến đảo chiều; tốc độ và động năng được bảo toàn. b) Sai: trái với hệ thức hoặc điều kiện đã nêu. c) Đúng: Từ $E=kT$, k có đơn vị J/K. d) Sai: trái với hệ thức hoặc điều kiện đã nêu
\item (Sai) Khi va chạm đàn hồi với thành bình đứng yên, đại lượng của phân tử đổi dấu là — phát biểu khối lượng là đúng.
\item (Đúng) Đơn vị của hằng số Boltzmann k là — phát biểu đúng là: J/K.
\item (Sai) Đơn vị của hằng số Boltzmann k là — phát biểu Pa/ $K$ là đúng.
\end{itemchoice}
}
\end{ex}

% ===== Câu 32 =====
% ID: L12C2B12-10-TL
% Mức: VD
\begin{ex}
% ID: L12C2B12-10-TL
% Mức: VD
Trong dạng “Công thức áp suất theo khối lượng riêng và tốc độ phân tử”, 4
\choiceTF

\loigiai{
$p\sim T$ khi n không đổi.
}
\end{ex}

% ===== Câu 33 =====
% ID: L12C2B12-10-TLN
% Mức: VD
\begin{ex}
% ID: L12C2B12-10-TLN
% Mức: VD
Trong dạng “Công thức áp suất theo khối lượng riêng và tốc độ phân tử”, Áp suất tăng 4 lần khi mật độ phân tử không đổi. Nhiệt độ tăng bao nhiêu lần?
\shortans{4}
\loigiai{
$p\sim T$ khi n không đổi.
}
\end{ex}

% ===== Câu 34 =====
% ID: L12C2B12-10-TN
% Mức: VDC
\dangbt{Bài toán thực tế về áp suất khí theo mô hình động học phân tử}
\begin{ex}
% ID: L12C2B12-10-TN
% Mức: VDC
Phát biểu nào phản ánh đúng bản chất bài toán?
\choice
{Có thể bỏ qua nhiệt độ tuyệt đối}
{Không cần xác định lượng khí}
{Mọi quá trình đều đẳng nhiệt}
{\True Phải dùng $p=\frac13\rho v_{rms}^2=nkT$ và kiểm tra khí loãng, cân bằng nhiệt và áp suất tuyệt đối}
\loigiai{
Đáp án D. Phải dùng $p=\frac13\rho v_{rms}^2=nkT$, đồng thời kiểm tra khí loãng, cân bằng nhiệt và áp suất tuyệt đối.
}
\end{ex}

% ===== Câu 35 =====
% ID: L12C2B12-11-DS
% Mức: TH
\dangbt{Liên hệ áp suất với động năng tịnh tiến trung bình}
\begin{ex}
% ID: L12C2B12-11-DS
% Mức: TH
Xét các phát biểu sau trong dạng “Liên hệ áp suất với động năng tịnh tiến trung bình”.
\choiceTF
{\True Áp suất chất khí lên thành bình có nguồn gốc từ — phát biểu đúng là: các va chạm của phân tử khí với thành bình.}
{Áp suất chất khí lên thành bình có nguồn gốc từ — phát biểu trọng lượng riêng của bình là đúng.}
{\True Công thức áp suất theo mô hình động học phân tử là — phát biểu đúng là: $p=\frac13\rho\overline{v^2}$.}
{Công thức áp suất theo mô hình động học phân tử là — phát biểu $p=\rho/3v^2$ là đúng.}
\loigiai{
\begin{itemchoice}
\item (Đúng) Áp suất chất khí lên thành bình có nguồn gốc từ — phát biểu đúng là: các va chạm của phân tử khí với thành bình. (Vì): Phân tử đổi động lượng khi va chạm và gây lực lên thành bình. b) Sai: trái với hệ thức hoặc điều kiện đã nêu. c) Đúng: Mô hình động học cho $p=\frac13\rho\overline{v^2}$. d) Sai: trái với hệ thức hoặc điều kiện đã nêu
\item (Sai) Áp suất chất khí lên thành bình có nguồn gốc từ — phát biểu trọng lượng riêng của bình là đúng.
\item (Đúng) Công thức áp suất theo mô hình động học phân tử là — phát biểu đúng là: $p=\frac13\rho\overline{v^2}$.
\item (Sai) Công thức áp suất theo mô hình động học phân tử là — phát biểu $p=\rho/3v^2$ là đúng.
\end{itemchoice}
}
\end{ex}

% ===== Câu 36 =====
% ID: L12C2B12-11-TL
% Mức: VD
\dangbt{Công thức áp suất theo khối lượng riêng và tốc độ phân tử}
\begin{ex}
% ID: L12C2B12-11-TL
% Mức: VD
Trong dạng “Công thức áp suất theo khối lượng riêng và tốc độ phân tử”, $v_{rms}\sim1/\sqrt m$, nên giảm 2 lần.
\choiceTF

\loigiai{
$v_{rms}\sim1/\sqrt m$, nên giảm 2 lần.
}
\end{ex}

% ===== Câu 37 =====
% ID: L12C2B12-11-TLN
% Mức: VD
\begin{ex}
% ID: L12C2B12-11-TLN
% Mức: VD
Trong dạng “Công thức áp suất theo khối lượng riêng và tốc độ phân tử”, Khối lượng phân tử tăng 4 lần ở cùng nhiệt độ. Tốc độ căn quân phương giảm bao nhiêu lần?
\shortans{2}
\loigiai{
$v_{rms}\sim1/\sqrt m$, nên giảm 2 lần.
}
\end{ex}

% ===== Câu 38 =====
% ID: L12C2B12-11-TN
% Mức: NB
\begin{ex}
% ID: L12C2B12-11-TN
% Mức: NB
Trong dạng “Công thức áp suất theo khối lượng riêng và tốc độ phân tử”, Áp suất chất khí lên thành bình có nguồn gốc từ
\choice
{\True các va chạm của phân tử khí với thành bình}
{trọng lượng riêng của bình}
{lực hút của Trái Đất}
{thể tích của phân tử}
\loigiai{
Áp suất chất khí lên thành bình có nguồn gốc từ các va chạm của các phân tử khí với thành bình. Khi các phân tử khí va chạm vào thành bình, chúng thay đổi động lượng. Theo định luật II Newton, sự thay đổi động lượng này gây ra một lực tác dụng lên thành bình. Lực này phân bố trên diện tích của thành bình, tạo nên áp suất
}
\end{ex}

% ===== Câu 39 =====
% ID: L12C2B12-12-DS
% Mức: TH
\dangbt{Liên hệ áp suất với động năng tịnh tiến trung bình}
\begin{ex}
% ID: L12C2B12-12-DS
% Mức: TH
Xét các phát biểu sau trong dạng “Liên hệ áp suất với động năng tịnh tiến trung bình”.
\choiceTF
{\True Khí có $\rho=1,5$ kg/m³ và tốc độ căn quân phương $500\,\text{m}$ /s. Áp suất gần bằng — phát biểu đúng là: $125\,\text{kPa}$.}
{Khí có $\rho=1,5$ kg/m³ và tốc độ căn quân phương $500\,\text{m}$ /s. Áp suất gần bằng — phát biểu $375\,\text{kPa}$ là đúng.}
{\True Khi nhiệt độ tuyệt đối tăng gấp 4, tốc độ căn quân phương của phân tử — phát biểu đúng là: tăng gấp 2.}
{Khi nhiệt độ tuyệt đối tăng gấp 4, tốc độ căn quân phương của phân tử — phát biểu giảm một nửa là đúng.}
\loigiai{
\begin{itemchoice}
\item (Đúng) Khí có $\rho=1,5$ kg/m³ và tốc độ căn quân phương $500\,\text{m}$ /s. Áp suất gần bằng — phát biểu đúng là: $125\,\text{kPa}$. (Vì): $p=\rho v_{rms}^2/3=125$ kPa. b) Sai: trái với hệ thức hoặc điều kiện đã nêu. c) Đúng: $v_{rms}\sim\sqrt{T}$. d) Sai: trái với hệ thức hoặc điều kiện đã nêu
\item (Sai) Khí có $\rho=1,5$ kg/m³ và tốc độ căn quân phương $500\,\text{m}$ /s. Áp suất gần bằng — phát biểu $375\,\text{kPa}$ là đúng.
\item (Đúng) Khi nhiệt độ tuyệt đối tăng gấp 4, tốc độ căn quân phương của phân tử — phát biểu đúng là: tăng gấp 2.
\item (Sai) Khi nhiệt độ tuyệt đối tăng gấp 4, tốc độ căn quân phương của phân tử — phát biểu giảm một nửa là đúng.
\end{itemchoice}
}
\end{ex}

% ===== Câu 40 =====
% ID: L12C2B12-12-TL
% Mức: VDC
\dangbt{Công thức áp suất theo khối lượng riêng và tốc độ phân tử}
\begin{bt}
% ID: L12C2B12-12-TL
% Mức: VDC
Trong dạng “Công thức áp suất theo khối lượng riêng và tốc độ phân tử”, undefined
\loigiai{
$\overline{E_k}=3kT/2$, nên $T$ tăng 3 lần.
}
\end{bt}

% ===== Câu 41 =====
% ID: L12C2B12-12-TLN
% Mức: VDC
\begin{ex}
% ID: L12C2B12-12-TLN
% Mức: VDC
Trong dạng “Công thức áp suất theo khối lượng riêng và tốc độ phân tử”, Động năng tịnh tiến trung bình tăng 3 lần. Nhiệt độ tuyệt đối tăng bao nhiêu lần?
\shortans{3}
\loigiai{
$\overline{E_k}=3kT/2$, nên $T$ tăng 3 lần.
}
\end{ex}

% ===== Câu 42 =====
% ID: L12C2B12-12-TN
% Mức: NB
\begin{ex}
% ID: L12C2B12-12-TN
% Mức: NB
Trong dạng “Công thức áp suất theo khối lượng riêng và tốc độ phân tử”, Công thức áp suất theo mô hình động học phân tử là
\choice
{\True $p=\frac13\rho\overline{v^2}$}
{$p=\rho gh$}
{$p=\rho/3v^2$}
{$p=3\rho\overline{v^2}$}
\loigiai{
Đáp án A. Mô hình động học cho $p=\frac13\rho\overline{v^2}$.
}
\end{ex}

% ===== Câu 43 =====
% ID: L12C2B12-13-DS
% Mức: VD
\dangbt{Liên hệ áp suất với động năng tịnh tiến trung bình}
\begin{ex}
% ID: L12C2B12-13-DS
% Mức: VD
Xét các phát biểu sau trong dạng “Liên hệ áp suất với động năng tịnh tiến trung bình”.
\choiceTF
{\True Động năng tịnh tiến trung bình của một phân tử khí lí tưởng bằng — phát biểu đúng là: $\frac32kT$.}
{Động năng tịnh tiến trung bình của một phân tử khí lí tưởng bằng — phát biểu $kT$ là đúng.}
{\True Hệ thức liên hệ áp suất với mật độ phân tử n là — phát biểu đúng là: $p=nkT$.}
{Hệ thức liên hệ áp suất với mật độ phân tử n là — phát biểu $p=k/(nT)$ là đúng.}
\loigiai{
\begin{itemchoice}
\item (Đúng) Động năng tịnh tiến trung bình của một phân tử khí lí tưởng bằng — phát biểu đúng là: $\frac32kT$. (Vì): Theo thuyết động học, $\overline{E_k}=3kT/2$. b) Sai: trái với hệ thức hoặc điều kiện đã nêu. c) Đúng: Với n là số phân tử trên một đơn vị thể tích, $p=nkT$. d) Sai: trái với hệ thức hoặc điều kiện đã nêu
\item (Sai) Động năng tịnh tiến trung bình của một phân tử khí lí tưởng bằng — phát biểu $kT$ là đúng.
\item (Đúng) Hệ thức liên hệ áp suất với mật độ phân tử n là — phát biểu đúng là: $p=nkT$.
\item (Sai) Hệ thức liên hệ áp suất với mật độ phân tử n là — phát biểu $p=k/(nT)$ là đúng.
\end{itemchoice}
}
\end{ex}

% ===== Câu 44 =====
% ID: L12C2B12-13-TL
% Mức: TH
\begin{ex}
% ID: L12C2B12-13-TL
% Mức: TH
Trong dạng “Liên hệ áp suất với động năng tịnh tiến trung bình”, Trình bày cơ sở lí thuyết của dạng “Liên hệ áp suất với động năng tịnh tiến trung bình” và nêu điều kiện áp dụng.
\choiceTF

\loigiai{
Nêu đúng đại lượng không đổi, hệ thức đặc trưng và yêu cầu dùng nhiệt độ kelvin, áp suất tuyệt đối khi có liên quan.
}
\end{ex}

% ===== Câu 45 =====
% ID: L12C2B12-13-TLN
% Mức: TH
\begin{ex}
% ID: L12C2B12-13-TLN
% Mức: TH
Trong dạng “Liên hệ áp suất với động năng tịnh tiến trung bình”, Khí có $\rho=1,2$ kg/m³, $v_{rms}=300$ m/s. Tính p theo kPa.
\shortans{36}
\loigiai{
$p=\rho v_{rms}^2/3=36$ kPa.
}
\end{ex}

% ===== Câu 46 =====
% ID: L12C2B12-13-TN
% Mức: TH
\dangbt{Công thức áp suất theo khối lượng riêng và tốc độ phân tử}
\begin{ex}
% ID: L12C2B12-13-TN
% Mức: TH
Trong dạng “Công thức áp suất theo khối lượng riêng và tốc độ phân tử”, Khí có $\rho=1,4$ kg/m³ và tốc độ căn quân phương $480\,\text{m}$ /s. Áp suất gần bằng
\choice
{\True $107,52\,\text{kPa}$}
{$322,56\,\text{kPa}$}
{$35,84\,\text{kPa}$}
{$0,67\,\text{kPa}$}
\loigiai{
Đáp án A. $p=\rho v_{rms}^2/3=107,52$ kPa.
}
\end{ex}

% ===== Câu 47 =====
% ID: L12C2B12-14-DS
% Mức: VD
\dangbt{Liên hệ áp suất với động năng tịnh tiến trung bình}
\begin{ex}
% ID: L12C2B12-14-DS
% Mức: VD
Xét các phát biểu sau trong dạng “Liên hệ áp suất với động năng tịnh tiến trung bình”.
\choiceTF
{\True Ở cùng nhiệt độ, phân tử khí có khối lượng nhỏ hơn có tốc độ căn quân phương — phát biểu đúng là: lớn hơn.}
{Ở cùng nhiệt độ, phân tử khí có khối lượng nhỏ hơn có tốc độ căn quân phương — phát biểu nhỏ hơn là đúng.}
{\True Nếu giữ mật độ phân tử không đổi và tăng $T$ gấp đôi thì áp suất — phát biểu đúng là: tăng gấp đôi.}
{Nếu giữ mật độ phân tử không đổi và tăng $T$ gấp đôi thì áp suất — phát biểu không đổi là đúng.}
\loigiai{
\begin{itemchoice}
\item (Đúng) Ở cùng nhiệt độ, phân tử khí có khối lượng nhỏ hơn có tốc độ căn quân phương — phát biểu đúng là: lớn hơn. (Vì): $v_{rms}=\sqrt{3kT/m}$. b) Sai: trái với hệ thức hoặc điều kiện đã nêu. c) Đúng: Từ $p=nkT$, p tỉ lệ thuận T. d) Sai: trái với hệ thức hoặc điều kiện đã nêu
\item (Sai) Ở cùng nhiệt độ, phân tử khí có khối lượng nhỏ hơn có tốc độ căn quân phương — phát biểu nhỏ hơn là đúng.
\item (Đúng) Nếu giữ mật độ phân tử không đổi và tăng $T$ gấp đôi thì áp suất — phát biểu đúng là: tăng gấp đôi.
\item (Sai) Nếu giữ mật độ phân tử không đổi và tăng $T$ gấp đôi thì áp suất — phát biểu không đổi là đúng.
\end{itemchoice}
}
\end{ex}

% ===== Câu 48 =====
% ID: L12C2B12-14-TL
% Mức: TH
\begin{ex}
% ID: L12C2B12-14-TL
% Mức: TH
Trong dạng “Liên hệ áp suất với động năng tịnh tiến trung bình”, Mô tả dạng đồ thị đặc trưng và cách nhận biết quá trình trong dạng “Liên hệ áp suất với động năng tịnh tiến trung bình”.
\choiceTF

\loigiai{
Xác định các trục, đại lượng không đổi và suy ra hình dạng đồ thị từ hệ thức vật lí tương ứng.
}
\end{ex}

% ===== Câu 49 =====
% ID: L12C2B12-14-TLN
% Mức: TH
\begin{ex}
% ID: L12C2B12-14-TLN
% Mức: TH
Trong dạng “Liên hệ áp suất với động năng tịnh tiến trung bình”, Nhiệt độ tuyệt đối tăng 9 lần. Tốc độ căn quân phương tăng bao nhiêu lần?
\shortans{3}
\loigiai{
$v_{rms}\sim\sqrt T$, nên tăng 3 lần.
}
\end{ex}

% ===== Câu 50 =====
% ID: L12C2B12-14-TN
% Mức: TH
\dangbt{Công thức áp suất theo khối lượng riêng và tốc độ phân tử}
\begin{ex}
% ID: L12C2B12-14-TN
% Mức: TH
Trong dạng “Công thức áp suất theo khối lượng riêng và tốc độ phân tử”, Khi nhiệt độ tuyệt đối tăng gấp 4, tốc độ căn quân phương của phân tử
\choice
{\True tăng gấp 2}
{tăng gấp 4}
{giảm một nửa}
{không đổi}
\loigiai{
Tốc độ căn quân phương của phân tử khí lí tưởng tỉ lệ thuận với căn bậc hai của nhiệt độ tuyệt đối của khối khí.
Mối quan hệ này được biểu diễn là $v_{rms} \sim \sqrt{T}$.
Khi nhiệt độ tuyệt đối $T$ tăng gấp 4 lần, tức là $T' = 4T$, thì tốc độ căn quân phương mới $v'_{rms}$ sẽ tỉ lệ với $\sqrt{T'} = \sqrt{4T} = 2\sqrt{T}$.
Vậy, tốc độ căn quân phương của phân tử tăng gấp 2 lần
}
\end{ex}

% ===== Câu 51 =====
% ID: L12C2B12-15-DS
% Mức: VDC
\dangbt{Liên hệ áp suất với động năng tịnh tiến trung bình}
\begin{ex}
% ID: L12C2B12-15-DS
% Mức: VDC
Xét các phát biểu sau trong dạng “Liên hệ áp suất với động năng tịnh tiến trung bình”.
\choiceTF
{\True Khi va chạm đàn hồi với thành bình đứng yên, đại lượng của phân tử đổi dấu là — phát biểu đúng là: thành phần vận tốc vuông góc thành.}
{Khi va chạm đàn hồi với thành bình đứng yên, đại lượng của phân tử đổi dấu là — phát biểu khối lượng là đúng.}
{\True Đơn vị của hằng số Boltzmann k là — phát biểu đúng là: J/K.}
{Đơn vị của hằng số Boltzmann k là — phát biểu Pa/ $K$ là đúng.}
\loigiai{
\begin{itemchoice}
\item (Đúng) Khi va chạm đàn hồi với thành bình đứng yên, đại lượng của phân tử đổi dấu là — phát biểu đúng là: thành phần vận tốc vuông góc thành. (Vì): Thành phần pháp tuyến đảo chiều; tốc độ và động năng được bảo toàn. b) Sai: trái với hệ thức hoặc điều kiện đã nêu. c) Đúng: Từ $E=kT$, k có đơn vị J/K. d) Sai: trái với hệ thức hoặc điều kiện đã nêu
\item (Sai) Khi va chạm đàn hồi với thành bình đứng yên, đại lượng của phân tử đổi dấu là — phát biểu khối lượng là đúng.
\item (Đúng) Đơn vị của hằng số Boltzmann k là — phát biểu đúng là: J/K.
\item (Sai) Đơn vị của hằng số Boltzmann k là — phát biểu Pa/ $K$ là đúng.
\end{itemchoice}
}
\end{ex}

% ===== Câu 52 =====
% ID: L12C2B12-15-TL
% Mức: VD
\begin{bt}
% ID: L12C2B12-15-TL
% Mức: VD
Trong dạng “Liên hệ áp suất với động năng tịnh tiến trung bình”, Mật độ phân tử tăng 4 lần, nhiệt độ không đổi. Áp suất tăng bao nhiêu lần?
\loigiai{
Từ $p=nkT$, p tăng 4 lần.
}
\end{bt}

% ===== Câu 53 =====
% ID: L12C2B12-15-TLN
% Mức: VD
\begin{ex}
% ID: L12C2B12-15-TLN
% Mức: VD
Trong dạng “Liên hệ áp suất với động năng tịnh tiến trung bình”, Mật độ phân tử tăng 4 lần, nhiệt độ không đổi. Áp suất tăng bao nhiêu lần?
\shortans{4}
\loigiai{
Từ $p=nkT$, p tăng 4 lần.
}
\end{ex}

% ===== Câu 54 =====
% ID: L12C2B12-15-TN
% Mức: TH
\dangbt{Công thức áp suất theo khối lượng riêng và tốc độ phân tử}
\begin{ex}
% ID: L12C2B12-15-TN
% Mức: TH
Trong dạng “Công thức áp suất theo khối lượng riêng và tốc độ phân tử”, Động năng tịnh tiến trung bình của một phân tử khí lí tưởng bằng
\choice
{\True $\frac32kT$}
{$kT$}
{$\frac12kT$}
{$3kT$}
\loigiai{
Đáp án A. Theo thuyết động học, $\overline{E_k}=3kT/2$.
}
\end{ex}

% ===== Câu 55 =====
% ID: L12C2B12-16-DS
% Mức: TH
\dangbt{Nguồn gốc áp suất khí và va chạm phân tử}
\begin{ex}
% ID: L12C2B12-16-DS
% Mức: TH
Xét các phát biểu sau trong dạng “Nguồn gốc áp suất khí và va chạm phân tử”.
\choiceTF
{\True Áp suất chất khí lên thành bình có nguồn gốc từ — phát biểu đúng là: các va chạm của phân tử khí với thành bình.}
{Áp suất chất khí lên thành bình có nguồn gốc từ — phát biểu trọng lượng riêng của bình là đúng.}
{\True Công thức áp suất theo mô hình động học phân tử là — phát biểu đúng là: $p=\frac13\rho\overline{v^2}$.}
{Công thức áp suất theo mô hình động học phân tử là — phát biểu $p=\rho/3v^2$ là đúng.}
\loigiai{
\begin{itemchoice}
\item (Đúng) Áp suất chất khí lên thành bình có nguồn gốc từ — phát biểu đúng là: các va chạm của phân tử khí với thành bình. (Vì): Phân tử đổi động lượng khi va chạm và gây lực lên thành bình. b) Sai: trái với hệ thức hoặc điều kiện đã nêu. c) Đúng: Mô hình động học cho $p=\frac13\rho\overline{v^2}$. d) Sai: trái với hệ thức hoặc điều kiện đã nêu
\item (Sai) Áp suất chất khí lên thành bình có nguồn gốc từ — phát biểu trọng lượng riêng của bình là đúng.
\item (Đúng) Công thức áp suất theo mô hình động học phân tử là — phát biểu đúng là: $p=\frac13\rho\overline{v^2}$.
\item (Sai) Công thức áp suất theo mô hình động học phân tử là — phát biểu $p=\rho/3v^2$ là đúng.
\end{itemchoice}
}
\end{ex}

% ===== Câu 56 =====
% ID: L12C2B12-16-TL
% Mức: VD
\dangbt{Liên hệ áp suất với động năng tịnh tiến trung bình}
\begin{ex}
% ID: L12C2B12-16-TL
% Mức: VD
Trong dạng “Liên hệ áp suất với động năng tịnh tiến trung bình”, 2
\choiceTF

\loigiai{
$p\sim T$ khi n không đổi.
}
\end{ex}

% ===== Câu 57 =====
% ID: L12C2B12-16-TLN
% Mức: VD
\begin{ex}
% ID: L12C2B12-16-TLN
% Mức: VD
Trong dạng “Liên hệ áp suất với động năng tịnh tiến trung bình”, Áp suất tăng 2 lần khi mật độ phân tử không đổi. Nhiệt độ tăng bao nhiêu lần?
\shortans{2}
\loigiai{
$p\sim T$ khi n không đổi.
}
\end{ex}

% ===== Câu 58 =====
% ID: L12C2B12-16-TN
% Mức: VD
\dangbt{Công thức áp suất theo khối lượng riêng và tốc độ phân tử}
\begin{ex}
% ID: L12C2B12-16-TN
% Mức: VD
Trong dạng “Công thức áp suất theo khối lượng riêng và tốc độ phân tử”, Hệ thức liên hệ áp suất với mật độ phân tử n là
\choice
{\True $p=nkT$}
{$p=nRT$}
{$p=k/(nT)$}
{$p=nk/T$}
\loigiai{
Đáp án A. Với n là số phân tử trên một đơn vị thể tích, $p=nkT$.
}
\end{ex}

% ===== Câu 59 =====
% ID: L12C2B12-17-DS
% Mức: TH
\dangbt{Nguồn gốc áp suất khí và va chạm phân tử}
\begin{ex}
% ID: L12C2B12-17-DS
% Mức: TH
Xét các phát biểu sau trong dạng “Nguồn gốc áp suất khí và va chạm phân tử”.
\choiceTF
{\True Khí có $\rho=1,3$ kg/m³ và tốc độ căn quân phương $460\,\text{m}$ /s. Áp suất gần bằng — phát biểu đúng là: $91,69\,\text{kPa}$.}
{Khí có $\rho=1,3$ kg/m³ và tốc độ căn quân phương $460\,\text{m}$ /s. Áp suất gần bằng — phát biểu $275,08\,\text{kPa}$ là đúng.}
{\True Khi nhiệt độ tuyệt đối tăng gấp 4, tốc độ căn quân phương của phân tử — phát biểu đúng là: tăng gấp 2.}
{Khi nhiệt độ tuyệt đối tăng gấp 4, tốc độ căn quân phương của phân tử — phát biểu giảm một nửa là đúng.}
\loigiai{
\begin{itemchoice}
\item (Đúng) Khí có $\rho=1,3$ kg/m³ và tốc độ căn quân phương $460\,\text{m}$ /s. Áp suất gần bằng — phát biểu đúng là: $91,69\,\text{kPa}$. (Vì): $p=\rho v_{rms}^2/3=91,69$ kPa. b) Sai: trái với hệ thức hoặc điều kiện đã nêu. c) Đúng: $v_{rms}\sim\sqrt{T}$. d) Sai: trái với hệ thức hoặc điều kiện đã nêu
\item (Sai) Khí có $\rho=1,3$ kg/m³ và tốc độ căn quân phương $460\,\text{m}$ /s. Áp suất gần bằng — phát biểu $275,08\,\text{kPa}$ là đúng.
\item (Đúng) Khi nhiệt độ tuyệt đối tăng gấp 4, tốc độ căn quân phương của phân tử — phát biểu đúng là: tăng gấp 2.
\item (Sai) Khi nhiệt độ tuyệt đối tăng gấp 4, tốc độ căn quân phương của phân tử — phát biểu giảm một nửa là đúng.
\end{itemchoice}
}
\end{ex}

% ===== Câu 60 =====
% ID: L12C2B12-17-TL
% Mức: VD
\dangbt{Liên hệ áp suất với động năng tịnh tiến trung bình}
\begin{ex}
% ID: L12C2B12-17-TL
% Mức: VD
Trong dạng “Liên hệ áp suất với động năng tịnh tiến trung bình”, $v_{rms}\sim1/\sqrt m$, nên giảm 3 lần.
\choiceTF

\loigiai{
$v_{rms}\sim1/\sqrt m$, nên giảm 3 lần.
}
\end{ex}

% ===== Câu 61 =====
% ID: L12C2B12-17-TLN
% Mức: VD
\begin{ex}
% ID: L12C2B12-17-TLN
% Mức: VD
Trong dạng “Liên hệ áp suất với động năng tịnh tiến trung bình”, Khối lượng phân tử tăng 9 lần ở cùng nhiệt độ. Tốc độ căn quân phương giảm bao nhiêu lần?
\shortans{3}
\loigiai{
$v_{rms}\sim1/\sqrt m$, nên giảm 3 lần.
}
\end{ex}

% ===== Câu 62 =====
% ID: L12C2B12-17-TN
% Mức: VD
\dangbt{Công thức áp suất theo khối lượng riêng và tốc độ phân tử}
\begin{ex}
% ID: L12C2B12-17-TN
% Mức: VD
Trong dạng “Công thức áp suất theo khối lượng riêng và tốc độ phân tử”, Ở cùng nhiệt độ, phân tử khí có khối lượng nhỏ hơn có tốc độ căn quân phương
\choice
{\True lớn hơn}
{nhỏ hơn}
{bằng nhau}
{bằng không}
\loigiai{
Đáp án A. $v_{rms}=\sqrt{3kT/m}$.
}
\end{ex}

% ===== Câu 63 =====
% ID: L12C2B12-18-DS
% Mức: VD
\dangbt{Nguồn gốc áp suất khí và va chạm phân tử}
\begin{ex}
% ID: L12C2B12-18-DS
% Mức: VD
Xét các phát biểu sau trong dạng “Nguồn gốc áp suất khí và va chạm phân tử”.
\choiceTF
{\True Động năng tịnh tiến trung bình của một phân tử khí lí tưởng bằng — phát biểu đúng là: $\frac32kT$.}
{Động năng tịnh tiến trung bình của một phân tử khí lí tưởng bằng — phát biểu $kT$ là đúng.}
{\True Hệ thức liên hệ áp suất với mật độ phân tử n là — phát biểu đúng là: $p=nkT$.}
{Hệ thức liên hệ áp suất với mật độ phân tử n là — phát biểu $p=k/(nT)$ là đúng.}
\loigiai{
\begin{itemchoice}
\item (Đúng) Động năng tịnh tiến trung bình của một phân tử khí lí tưởng bằng — phát biểu đúng là: $\frac32kT$. (Vì): Theo thuyết động học, $\overline{E_k}=3kT/2$. b) Sai: trái với hệ thức hoặc điều kiện đã nêu. c) Đúng: Với n là số phân tử trên một đơn vị thể tích, $p=nkT$. d) Sai: trái với hệ thức hoặc điều kiện đã nêu
\item (Sai) Động năng tịnh tiến trung bình của một phân tử khí lí tưởng bằng — phát biểu $kT$ là đúng.
\item (Đúng) Hệ thức liên hệ áp suất với mật độ phân tử n là — phát biểu đúng là: $p=nkT$.
\item (Sai) Hệ thức liên hệ áp suất với mật độ phân tử n là — phát biểu $p=k/(nT)$ là đúng.
\end{itemchoice}
}
\end{ex}

% ===== Câu 64 =====
% ID: L12C2B12-18-TL
% Mức: VDC
\dangbt{Liên hệ áp suất với động năng tịnh tiến trung bình}
\begin{bt}
% ID: L12C2B12-18-TL
% Mức: VDC
Trong dạng “Liên hệ áp suất với động năng tịnh tiến trung bình”, undefined
\loigiai{
$\overline{E_k}=3kT/2$, nên $T$ tăng 4 lần.
}
\end{bt}

% ===== Câu 65 =====
% ID: L12C2B12-18-TLN
% Mức: VDC
\begin{ex}
% ID: L12C2B12-18-TLN
% Mức: VDC
Trong dạng “Liên hệ áp suất với động năng tịnh tiến trung bình”, Động năng tịnh tiến trung bình tăng 4 lần. Nhiệt độ tuyệt đối tăng bao nhiêu lần?
\shortans{4}
\loigiai{
$\overline{E_k}=3kT/2$, nên $T$ tăng 4 lần.
}
\end{ex}

% ===== Câu 66 =====
% ID: L12C2B12-18-TN
% Mức: VD
\dangbt{Công thức áp suất theo khối lượng riêng và tốc độ phân tử}
\begin{ex}
% ID: L12C2B12-18-TN
% Mức: VD
Trong dạng “Công thức áp suất theo khối lượng riêng và tốc độ phân tử”, Nếu giữ mật độ phân tử không đổi và tăng $T$ gấp đôi thì áp suất
\choice
{\True tăng gấp đôi}
{giảm một nửa}
{không đổi}
{tăng gấp bốn}
\loigiai{
Đáp án A. Từ $p=nkT$, p tỉ lệ thuận T.
}
\end{ex}

% ===== Câu 67 =====
% ID: L12C2B12-19-DS
% Mức: VD
\dangbt{Nguồn gốc áp suất khí và va chạm phân tử}
\begin{ex}
% ID: L12C2B12-19-DS
% Mức: VD
Xét các phát biểu sau trong dạng “Nguồn gốc áp suất khí và va chạm phân tử”.
\choiceTF
{\True Ở cùng nhiệt độ, phân tử khí có khối lượng nhỏ hơn có tốc độ căn quân phương — phát biểu đúng là: lớn hơn.}
{Ở cùng nhiệt độ, phân tử khí có khối lượng nhỏ hơn có tốc độ căn quân phương — phát biểu nhỏ hơn là đúng.}
{\True Nếu giữ mật độ phân tử không đổi và tăng $T$ gấp đôi thì áp suất — phát biểu đúng là: tăng gấp đôi.}
{Nếu giữ mật độ phân tử không đổi và tăng $T$ gấp đôi thì áp suất — phát biểu không đổi là đúng.}
\loigiai{
\begin{itemchoice}
\item (Đúng) Ở cùng nhiệt độ, phân tử khí có khối lượng nhỏ hơn có tốc độ căn quân phương — phát biểu đúng là: lớn hơn. (Vì): $v_{rms}=\sqrt{3kT/m}$. b) Sai: trái với hệ thức hoặc điều kiện đã nêu. c) Đúng: Từ $p=nkT$, p tỉ lệ thuận T. d) Sai: trái với hệ thức hoặc điều kiện đã nêu
\item (Sai) Ở cùng nhiệt độ, phân tử khí có khối lượng nhỏ hơn có tốc độ căn quân phương — phát biểu nhỏ hơn là đúng.
\item (Đúng) Nếu giữ mật độ phân tử không đổi và tăng $T$ gấp đôi thì áp suất — phát biểu đúng là: tăng gấp đôi.
\item (Sai) Nếu giữ mật độ phân tử không đổi và tăng $T$ gấp đôi thì áp suất — phát biểu không đổi là đúng.
\end{itemchoice}
}
\end{ex}

% ===== Câu 68 =====
% ID: L12C2B12-19-TL
% Mức: TH
\begin{ex}
% ID: L12C2B12-19-TL
% Mức: TH
Trong dạng “Nguồn gốc áp suất khí và va chạm phân tử”, Trình bày cơ sở lí thuyết của dạng “Nguồn gốc áp suất khí và va chạm phân tử” và nêu điều kiện áp dụng.
\choiceTF

\loigiai{
Nêu đúng đại lượng không đổi, hệ thức đặc trưng và yêu cầu dùng nhiệt độ kelvin, áp suất tuyệt đối khi có liên quan.
}
\end{ex}

% ===== Câu 69 =====
% ID: L12C2B12-19-TLN
% Mức: TH
\begin{ex}
% ID: L12C2B12-19-TLN
% Mức: TH
Trong dạng “Nguồn gốc áp suất khí và va chạm phân tử”, Khí có $\rho=1,2$ kg/m³, $v_{rms}=300$ m/s. Tính p theo kPa.
\shortans{36}
\loigiai{
$p=\rho v_{rms}^2/3=36$ kPa.
}
\end{ex}

% ===== Câu 70 =====
% ID: L12C2B12-19-TN
% Mức: VDC
\dangbt{Công thức áp suất theo khối lượng riêng và tốc độ phân tử}
\begin{ex}
% ID: L12C2B12-19-TN
% Mức: VDC
Trong dạng “Công thức áp suất theo khối lượng riêng và tốc độ phân tử”, Khi va chạm đàn hồi với thành bình đứng yên, đại lượng của phân tử đổi dấu là
\choice
{\True thành phần vận tốc vuông góc thành}
{khối lượng}
{động năng}
{độ lớn vận tốc}
\loigiai{
Đáp án A. Thành phần pháp tuyến đảo chiều; tốc độ và động năng được bảo toàn.
}
\end{ex}

% ===== Câu 71 =====
% ID: L12C2B12-20-DS
% Mức: VDC
\dangbt{Nguồn gốc áp suất khí và va chạm phân tử}
\begin{ex}
% ID: L12C2B12-20-DS
% Mức: VDC
Xét các phát biểu sau trong dạng “Nguồn gốc áp suất khí và va chạm phân tử”.
\choiceTF
{\True Khi va chạm đàn hồi với thành bình đứng yên, đại lượng của phân tử đổi dấu là — phát biểu đúng là: thành phần vận tốc vuông góc thành.}
{Khi va chạm đàn hồi với thành bình đứng yên, đại lượng của phân tử đổi dấu là — phát biểu khối lượng là đúng.}
{\True Đơn vị của hằng số Boltzmann k là — phát biểu đúng là: J/K.}
{Đơn vị của hằng số Boltzmann k là — phát biểu Pa/ $K$ là đúng.}
\loigiai{
\begin{itemchoice}
\item (Đúng) Khi va chạm đàn hồi với thành bình đứng yên, đại lượng của phân tử đổi dấu là — phát biểu đúng là: thành phần vận tốc vuông góc thành. (Vì): Thành phần pháp tuyến đảo chiều; tốc độ và động năng được bảo toàn. b) Sai: trái với hệ thức hoặc điều kiện đã nêu. c) Đúng: Từ $E=kT$, k có đơn vị J/K. d) Sai: trái với hệ thức hoặc điều kiện đã nêu
\item (Sai) Khi va chạm đàn hồi với thành bình đứng yên, đại lượng của phân tử đổi dấu là — phát biểu khối lượng là đúng.
\item (Đúng) Đơn vị của hằng số Boltzmann k là — phát biểu đúng là: J/K.
\item (Sai) Đơn vị của hằng số Boltzmann k là — phát biểu Pa/ $K$ là đúng.
\end{itemchoice}
}
\end{ex}

% ===== Câu 72 =====
% ID: L12C2B12-20-TL
% Mức: TH
\begin{ex}
% ID: L12C2B12-20-TL
% Mức: TH
Trong dạng “Nguồn gốc áp suất khí và va chạm phân tử”, Mô tả dạng đồ thị đặc trưng và cách nhận biết quá trình trong dạng “Nguồn gốc áp suất khí và va chạm phân tử”.
\choiceTF

\loigiai{
Xác định các trục, đại lượng không đổi và suy ra hình dạng đồ thị từ hệ thức vật lí tương ứng.
}
\end{ex}

% ===== Câu 73 =====
% ID: L12C2B12-20-TLN
% Mức: TH
\begin{ex}
% ID: L12C2B12-20-TLN
% Mức: TH
Trong dạng “Nguồn gốc áp suất khí và va chạm phân tử”, Nhiệt độ tuyệt đối tăng 16 lần. Tốc độ căn quân phương tăng bao nhiêu lần?
\shortans{4}
\loigiai{
$v_{rms}\sim\sqrt T$, nên tăng 4 lần.
}
\end{ex}

% ===== Câu 74 =====
% ID: L12C2B12-20-TN
% Mức: VDC
\dangbt{Công thức áp suất theo khối lượng riêng và tốc độ phân tử}
\begin{ex}
% ID: L12C2B12-20-TN
% Mức: VDC
Trong dạng “Công thức áp suất theo khối lượng riêng và tốc độ phân tử”, Đơn vị của hằng số Boltzmann k là
\choice
{\True J/ $K$}
{J· $K$}
{Pa/ $K$}
{mol/ $J$}
\loigiai{
Đáp án A. Từ $E=kT$, k có đơn vị J/K.
}
\end{ex}

% ===== Câu 75 =====
% ID: L12C2B12-21-DS
% Mức: TH
\dangbt{So sánh chuyển động nhiệt của các loại khí}
\begin{ex}
% ID: L12C2B12-21-DS
% Mức: TH
Xét dạng “So sánh chuyển động nhiệt của các loại khí”: a) Công thức đặc trưng là $v_{rms}=\sqrt{3RT/M}$. b) Có thể bỏ qua điều kiện các khí lí tưởng được so sánh ở nhiệt độ xác định. c) ở cùng $T$, $v_{rms}\sim1/\sqrt M$. d) Có thể cho rằng mọi loại phân tử ở cùng nhiệt độ có cùng tốc độ.
\choiceTF

\loigiai{
a) Đúng: phù hợp với công thức và điều kiện đã nêu. b) Sai: phát biểu không thỏa công thức hoặc điều kiện áp dụng. c) Đúng: phù hợp với công thức và điều kiện đã nêu. d) Sai: phát biểu không thỏa công thức hoặc điều kiện áp dụng.
}
\end{ex}

% ===== Câu 76 =====
% ID: L12C2B12-21-TLN
% Mức: VD
\dangbt{Nguồn gốc áp suất khí và va chạm phân tử}
\begin{ex}
% ID: L12C2B12-21-TLN
% Mức: VD
Trong dạng “Nguồn gốc áp suất khí và va chạm phân tử”, Mật độ phân tử tăng 2 lần, nhiệt độ không đổi. Áp suất tăng bao nhiêu lần?
\shortans{2}
\loigiai{
Từ $p=nkT$, p tăng 2 lần.
}
\end{ex}

% ===== Câu 77 =====
% ID: L12C2B12-21-TN
% Mức: NB
\dangbt{Liên hệ áp suất với động năng tịnh tiến trung bình}
\begin{ex}
% ID: L12C2B12-21-TN
% Mức: NB
Trong dạng “Liên hệ áp suất với động năng tịnh tiến trung bình”, Áp suất chất khí lên thành bình có nguồn gốc từ
\choice
{\True các va chạm của phân tử khí với thành bình}
{trọng lượng riêng của bình}
{lực hút của Trái Đất בלבד}
{thể tích của phân tử}
\loigiai{
Đáp án A. Phân tử đổi động lượng khi va chạm và gây lực lên thành bình.
}
\end{ex}

% ===== Câu 78 =====
% ID: L12C2B12-22-DS
% Mức: TH
\dangbt{So sánh chuyển động nhiệt của các loại khí}
\begin{ex}
% ID: L12C2B12-22-DS
% Mức: TH
Xét dạng “So sánh chuyển động nhiệt của các loại khí”: a) Cần dùng đơn vị phù hợp: m/s. b) Nhiệt độ tuyệt đối có thể mang giá trị âm. c) Phải xác định đại lượng không đổi trước khi tính. d) Không cần kiểm tra thứ nguyên.
\choiceTF

\loigiai{
a) Đúng: phù hợp với công thức và điều kiện đã nêu. b) Sai: phát biểu không thỏa công thức hoặc điều kiện áp dụng. c) Đúng: phù hợp với công thức và điều kiện đã nêu. d) Sai: phát biểu không thỏa công thức hoặc điều kiện áp dụng.
}
\end{ex}

% ===== Câu 79 =====
% ID: L12C2B12-22-TLN
% Mức: VD
\dangbt{Nguồn gốc áp suất khí và va chạm phân tử}
\begin{ex}
% ID: L12C2B12-22-TLN
% Mức: VD
Trong dạng “Nguồn gốc áp suất khí và va chạm phân tử”, Áp suất tăng 3 lần khi mật độ phân tử không đổi. Nhiệt độ tăng bao nhiêu lần?
\shortans{3}
\loigiai{
$p\sim T$ khi n không đổi.
}
\end{ex}

% ===== Câu 80 =====
% ID: L12C2B12-22-TN
% Mức: NB
\dangbt{Liên hệ áp suất với động năng tịnh tiến trung bình}
\begin{ex}
% ID: L12C2B12-22-TN
% Mức: NB
Trong dạng “Liên hệ áp suất với động năng tịnh tiến trung bình”, Công thức áp suất theo mô hình động học phân tử là
\choice
{\True $p=\frac13\rho\overline{v^2}$}
{$p=\rho gh$}
{$p=\rho/3v^2$}
{$p=3\rho\overline{v^2}$}
\loigiai{
Đáp án A. Mô hình động học cho $p=\frac13\rho\overline{v^2}$.
}
\end{ex}

% ===== Câu 81 =====
% ID: L12C2B12-23-DS
% Mức: VD
\dangbt{So sánh chuyển động nhiệt của các loại khí}
\begin{ex}
% ID: L12C2B12-23-DS
% Mức: VD
Xét dạng “So sánh chuyển động nhiệt của các loại khí”: a) Bài mẫu cho kết quả 2. b) Kết quả bài mẫu gấp đôi 2. c) Lời giải phải nêu công thức trước khi thay số. d) Mọi dữ kiện số đều không cần đổi đơn vị.
\choiceTF

\loigiai{
a) Đúng: phù hợp với công thức và điều kiện đã nêu. b) Sai: phát biểu không thỏa công thức hoặc điều kiện áp dụng. c) Đúng: phù hợp với công thức và điều kiện đã nêu. d) Sai: phát biểu không thỏa công thức hoặc điều kiện áp dụng.
}
\end{ex}

% ===== Câu 82 =====
% ID: L12C2B12-23-TLN
% Mức: VD
\dangbt{Nguồn gốc áp suất khí và va chạm phân tử}
\begin{ex}
% ID: L12C2B12-23-TLN
% Mức: VD
Trong dạng “Nguồn gốc áp suất khí và va chạm phân tử”, Khối lượng phân tử tăng 16 lần ở cùng nhiệt độ. Tốc độ căn quân phương giảm bao nhiêu lần?
\shortans{4}
\loigiai{
$v_{rms}\sim1/\sqrt m$, nên giảm 4 lần.
}
\end{ex}

% ===== Câu 83 =====
% ID: L12C2B12-23-TN
% Mức: TH
\dangbt{Liên hệ áp suất với động năng tịnh tiến trung bình}
\begin{ex}
% ID: L12C2B12-23-TN
% Mức: TH
Trong dạng “Liên hệ áp suất với động năng tịnh tiến trung bình”, Khí có $\rho=1,5$ kg/m³ và tốc độ căn quân phương $500\,\text{m}$ /s. Áp suất gần bằng
\choice
{\True $125\,\text{kPa}$}
{$375\,\text{kPa}$}
{$41,67\,\text{kPa}$}
{$0,75\,\text{kPa}$}
\loigiai{
Đáp án A. $p=\rho v_{rms}^2/3=125$ kPa.
}
\end{ex}

% ===== Câu 84 =====
% ID: L12C2B12-24-DS
% Mức: VD
\dangbt{So sánh chuyển động nhiệt của các loại khí}
\begin{ex}
% ID: L12C2B12-24-DS
% Mức: VD
Xét dạng “So sánh chuyển động nhiệt của các loại khí”: a) Quan hệ cần dùng là ở cùng $T$, $v_{rms}\sim1/\sqrt M$. b) Có thể suy ra quan hệ mà không xét điều kiện quá trình. c) Kết quả phải phù hợp chiều biến thiên vật lí. d) Áp suất dư luôn thay được áp suất tuyệt đối.
\choiceTF

\loigiai{
a) Đúng: phù hợp với công thức và điều kiện đã nêu. b) Sai: phát biểu không thỏa công thức hoặc điều kiện áp dụng. c) Đúng: phù hợp với công thức và điều kiện đã nêu. d) Sai: phát biểu không thỏa công thức hoặc điều kiện áp dụng.
}
\end{ex}

% ===== Câu 85 =====
% ID: L12C2B12-24-TLN
% Mức: VDC
\dangbt{Nguồn gốc áp suất khí và va chạm phân tử}
\begin{ex}
% ID: L12C2B12-24-TLN
% Mức: VDC
Trong dạng “Nguồn gốc áp suất khí và va chạm phân tử”, Động năng tịnh tiến trung bình tăng 2 lần. Nhiệt độ tuyệt đối tăng bao nhiêu lần?
\shortans{2}
\loigiai{
$\overline{E_k}=3kT/2$, nên $T$ tăng 2 lần.
}
\end{ex}

% ===== Câu 86 =====
% ID: L12C2B12-24-TN
% Mức: TH
\dangbt{Liên hệ áp suất với động năng tịnh tiến trung bình}
\begin{ex}
% ID: L12C2B12-24-TN
% Mức: TH
Trong dạng “Liên hệ áp suất với động năng tịnh tiến trung bình”, Khi nhiệt độ tuyệt đối tăng gấp 4, tốc độ căn quân phương của phân tử
\choice
{\True tăng gấp 2}
{tăng gấp 4}
{giảm một nửa}
{không đổi}
\loigiai{
Đáp án A. $v_{rms}\sim\sqrt{T}$.
}
\end{ex}

% ===== Câu 87 =====
% ID: L12C2B12-25-DS
% Mức: VDC
\dangbt{So sánh chuyển động nhiệt của các loại khí}
\begin{ex}
% ID: L12C2B12-25-DS
% Mức: VDC
Xét dạng “So sánh chuyển động nhiệt của các loại khí”: a) Dạng này có thể yêu cầu phối hợp đổi đơn vị và lập tỉ số. b) Kelvin được đổi sang Celsius bằng cách cộng 273. c) Cần ghi rõ đại lượng cần tìm. d) Công thức nào có đủ kí hiệu cũng dùng được.
\choiceTF

\loigiai{
a) Đúng: phù hợp với công thức và điều kiện đã nêu. b) Sai: phát biểu không thỏa công thức hoặc điều kiện áp dụng. c) Đúng: phù hợp với công thức và điều kiện đã nêu. d) Sai: phát biểu không thỏa công thức hoặc điều kiện áp dụng.
}
\end{ex}

% ===== Câu 88 =====
% ID: L12C2B12-25-TLN
% Mức: TH
\begin{ex}
% ID: L12C2B12-25-TLN
% Mức: TH
Ở cùng nhiệt độ, khối lượng mol của khí $B$ gấp 4 lần khí A. Tính tỉ số $v_A/v_B$.
\shortans{2}
\loigiai{
Ở cùng $T$, $v\sim1/\sqrt M$, nên $v_A/v_B=\sqrt{M_B/M_A}=2$.
}
\end{ex}

% ===== Câu 89 =====
% ID: L12C2B12-25-TN
% Mức: TH
\dangbt{Liên hệ áp suất với động năng tịnh tiến trung bình}
\begin{ex}
% ID: L12C2B12-25-TN
% Mức: TH
Trong dạng “Liên hệ áp suất với động năng tịnh tiến trung bình”, Động năng tịnh tiến trung bình của một phân tử khí lí tưởng bằng
\choice
{\True $\frac32kT$}
{$kT$}
{$\frac12kT$}
{$3kT$}
\loigiai{
Đáp án A. Theo thuyết động học, $\overline{E_k}=3kT/2$.
}
\end{ex}

% ===== Câu 90 =====
% ID: L12C2B12-26-DS
% Mức: TH
\dangbt{Tính tốc độ căn quân phương của phân tử khí}
\begin{ex}
% ID: L12C2B12-26-DS
% Mức: TH
Xét dạng “Tính tốc độ căn quân phương của phân tử khí”: a) Công thức đặc trưng là $v_{rms}=\sqrt{3RT/M}=\sqrt{3kT/m}$. b) Có thể bỏ qua điều kiện khí lí tưởng ở nhiệt độ tuyệt đối T. c) $v_{rms}$ tỉ lệ với $\sqrt{T/M}$. d) Có thể dùng khối lượng mol theo g/mol mà chưa đổi sang kg/mol.
\choiceTF

\loigiai{
a) Đúng: phù hợp với công thức và điều kiện đã nêu. b) Sai: phát biểu không thỏa công thức hoặc điều kiện áp dụng. c) Đúng: phù hợp với công thức và điều kiện đã nêu. d) Sai: phát biểu không thỏa công thức hoặc điều kiện áp dụng.
}
\end{ex}

% ===== Câu 91 =====
% ID: L12C2B12-26-TLN
% Mức: TH
\dangbt{So sánh chuyển động nhiệt của các loại khí}
\begin{ex}
% ID: L12C2B12-26-TLN
% Mức: TH
Đổi 27 ° $C$ sang nhiệt độ tuyệt đối theo kelvin.
\shortans{300}
\loigiai{
$T=27+273=300$ K.
}
\end{ex}

% ===== Câu 92 =====
% ID: L12C2B12-26-TN
% Mức: VD
\dangbt{Liên hệ áp suất với động năng tịnh tiến trung bình}
\begin{ex}
% ID: L12C2B12-26-TN
% Mức: VD
Trong dạng “Liên hệ áp suất với động năng tịnh tiến trung bình”, Hệ thức liên hệ áp suất với mật độ phân tử n là
\choice
{\True $p=nkT$}
{$p=nRT$}
{$p=k/(nT)$}
{$p=nk/T$}
\loigiai{
Đáp án A. Với n là số phân tử trên một đơn vị thể tích, $p=nkT$.
}
\end{ex}

% ===== Câu 93 =====
% ID: L12C2B12-27-DS
% Mức: TH
\dangbt{Tính tốc độ căn quân phương của phân tử khí}
\begin{ex}
% ID: L12C2B12-27-DS
% Mức: TH
Xét dạng “Tính tốc độ căn quân phương của phân tử khí”: a) Cần dùng đơn vị phù hợp: m/s. b) Nhiệt độ tuyệt đối có thể mang giá trị âm. c) Phải xác định đại lượng không đổi trước khi tính. d) Không cần kiểm tra thứ nguyên.
\choiceTF

\loigiai{
a) Đúng: phù hợp với công thức và điều kiện đã nêu. b) Sai: phát biểu không thỏa công thức hoặc điều kiện áp dụng. c) Đúng: phù hợp với công thức và điều kiện đã nêu. d) Sai: phát biểu không thỏa công thức hoặc điều kiện áp dụng.
}
\end{ex}

% ===== Câu 94 =====
% ID: L12C2B12-27-TLN
% Mức: VD
\dangbt{So sánh chuyển động nhiệt của các loại khí}
\begin{ex}
% ID: L12C2B12-27-TLN
% Mức: VD
Một đại lượng tăng từ 100 lên 150. Tính tỉ số giá trị sau so với trước.
\shortans{1,5}
\loigiai{
Tỉ số bằng $150/100=1,5$.
}
\end{ex}

% ===== Câu 95 =====
% ID: L12C2B12-27-TN
% Mức: VD
\dangbt{Liên hệ áp suất với động năng tịnh tiến trung bình}
\begin{ex}
% ID: L12C2B12-27-TN
% Mức: VD
Trong dạng “Liên hệ áp suất với động năng tịnh tiến trung bình”, Ở cùng nhiệt độ, phân tử khí có khối lượng nhỏ hơn có tốc độ căn quân phương
\choice
{\True lớn hơn}
{nhỏ hơn}
{bằng nhau}
{bằng không}
\loigiai{
Đáp án A. $v_{rms}=\sqrt{3kT/m}$.
}
\end{ex}

% ===== Câu 96 =====
% ID: L12C2B12-28-DS
% Mức: VD
\dangbt{Tính tốc độ căn quân phương của phân tử khí}
\begin{ex}
% ID: L12C2B12-28-DS
% Mức: VD
Xét dạng “Tính tốc độ căn quân phương của phân tử khí”: a) Bài mẫu cho kết quả 517. b) Kết quả bài mẫu gấp đôi 517. c) Lời giải phải nêu công thức trước khi thay số. d) Mọi dữ kiện số đều không cần đổi đơn vị.
\choiceTF

\loigiai{
a) Đúng: phù hợp với công thức và điều kiện đã nêu. b) Sai: phát biểu không thỏa công thức hoặc điều kiện áp dụng. c) Đúng: phù hợp với công thức và điều kiện đã nêu. d) Sai: phát biểu không thỏa công thức hoặc điều kiện áp dụng.
}
\end{ex}

% ===== Câu 97 =====
% ID: L12C2B12-28-TLN
% Mức: VD
\dangbt{So sánh chuyển động nhiệt của các loại khí}
\begin{ex}
% ID: L12C2B12-28-TLN
% Mức: VD
Một thể tích 2,5 $L$ bằng bao nhiêu m³? Viết kết quả theo dạng $a\cdot10^{-3}$.
\shortans{2,5}
\loigiai{
$2,5L=2,5\cdot10^{-3}$ m³.
}
\end{ex}

% ===== Câu 98 =====
% ID: L12C2B12-28-TN
% Mức: VD
\dangbt{Liên hệ áp suất với động năng tịnh tiến trung bình}
\begin{ex}
% ID: L12C2B12-28-TN
% Mức: VD
Trong dạng “Liên hệ áp suất với động năng tịnh tiến trung bình”, Nếu giữ mật độ phân tử không đổi và tăng $T$ gấp đôi thì áp suất
\choice
{\True tăng gấp đôi}
{giảm một nửa}
{không đổi}
{tăng gấp bốn}
\loigiai{
Đáp án A. Từ $p=nkT$, p tỉ lệ thuận T.
}
\end{ex}

% ===== Câu 99 =====
% ID: L12C2B12-29-DS
% Mức: VD
\dangbt{Tính tốc độ căn quân phương của phân tử khí}
\begin{ex}
% ID: L12C2B12-29-DS
% Mức: VD
Xét dạng “Tính tốc độ căn quân phương của phân tử khí”: a) Quan hệ cần dùng là $v_{rms}$ tỉ lệ với $\sqrt{T/M}$. b) Có thể suy ra quan hệ mà không xét điều kiện quá trình. c) Kết quả phải phù hợp chiều biến thiên vật lí. d) Áp suất dư luôn thay được áp suất tuyệt đối.
\choiceTF

\loigiai{
a) Đúng: phù hợp với công thức và điều kiện đã nêu. b) Sai: phát biểu không thỏa công thức hoặc điều kiện áp dụng. c) Đúng: phù hợp với công thức và điều kiện đã nêu. d) Sai: phát biểu không thỏa công thức hoặc điều kiện áp dụng.
}
\end{ex}

% ===== Câu 100 =====
% ID: L12C2B12-29-TLN
% Mức: VD
\dangbt{So sánh chuyển động nhiệt của các loại khí}
\begin{ex}
% ID: L12C2B12-29-TLN
% Mức: VD
Áp suất $120\,\text{kPa}$ bằng bao nhiêu Pa, lấy hệ số của $10^3$ làm đáp án?
\shortans{120}
\loigiai{
$120\,\text{kPa}=120\cdot10^3$ Pa.
}
\end{ex}

% ===== Câu 101 =====
% ID: L12C2B12-29-TN
% Mức: VDC
\dangbt{Liên hệ áp suất với động năng tịnh tiến trung bình}
\begin{ex}
% ID: L12C2B12-29-TN
% Mức: VDC
Trong dạng “Liên hệ áp suất với động năng tịnh tiến trung bình”, Khi va chạm đàn hồi với thành bình đứng yên, đại lượng của phân tử đổi dấu là
\choice
{\True thành phần vận tốc vuông góc thành}
{khối lượng}
{động năng}
{độ lớn vận tốc}
\loigiai{
Đáp án A. Thành phần pháp tuyến đảo chiều; tốc độ và động năng được bảo toàn.
}
\end{ex}

% ===== Câu 102 =====
% ID: L12C2B12-30-DS
% Mức: VDC
\dangbt{Tính tốc độ căn quân phương của phân tử khí}
\begin{ex}
% ID: L12C2B12-30-DS
% Mức: VDC
Xét dạng “Tính tốc độ căn quân phương của phân tử khí”: a) Dạng này có thể yêu cầu phối hợp đổi đơn vị và lập tỉ số. b) Kelvin được đổi sang Celsius bằng cách cộng 273. c) Cần ghi rõ đại lượng cần tìm. d) Công thức nào có đủ kí hiệu cũng dùng được.
\choiceTF

\loigiai{
a) Đúng: phù hợp với công thức và điều kiện đã nêu. b) Sai: phát biểu không thỏa công thức hoặc điều kiện áp dụng. c) Đúng: phù hợp với công thức và điều kiện đã nêu. d) Sai: phát biểu không thỏa công thức hoặc điều kiện áp dụng.
}
\end{ex}

% ===== Câu 103 =====
% ID: L12C2B12-30-TLN
% Mức: VDC
\dangbt{So sánh chuyển động nhiệt của các loại khí}
\begin{ex}
% ID: L12C2B12-30-TLN
% Mức: VDC
Theo quan hệ “ở cùng $T$, $v_{rms}\sim1/\sqrt M$ ”, nếu đại lượng phụ thuộc tỉ lệ thuận tăng gấp 4 thì hệ số tăng là bao nhiêu?
\shortans{4}
\loigiai{
Với quan hệ tỉ lệ thuận, hệ số biến đổi được giữ nguyên: 4.
}
\end{ex}

% ===== Câu 104 =====
% ID: L12C2B12-30-TN
% Mức: VDC
\dangbt{Liên hệ áp suất với động năng tịnh tiến trung bình}
\begin{ex}
% ID: L12C2B12-30-TN
% Mức: VDC
Trong dạng “Liên hệ áp suất với động năng tịnh tiến trung bình”, Đơn vị của hằng số Boltzmann k là
\choice
{\True J/ $K$}
{J· $K$}
{Pa/ $K$}
{mol/ $J$}
\loigiai{
Đáp án A. Từ $E=kT$, k có đơn vị J/K.
}
\end{ex}

% ===== Câu 105 =====
% ID: L12C2B12-31-TL
% Mức: TH
\dangbt{Tính tốc độ căn quân phương của phân tử khí}
\begin{ex}
% ID: L12C2B12-31-TL
% Mức: TH
Trình bày cơ sở lí thuyết và điều kiện áp dụng của dạng “Tính tốc độ căn quân phương của phân tử khí”.
\choiceTF

\loigiai{
$v_{rms}=\sqrt{3RT/M}=\sqrt{3kT/m}$. Điều kiện: khí lí tưởng ở nhiệt độ tuyệt đối T. Chuẩn hóa đơn vị m/s; sau đó lập phương trình và kiểm tra kết quả.
}
\end{ex}

% ===== Câu 106 =====
% ID: L12C2B12-31-TLN
% Mức: TH
\begin{ex}
% ID: L12C2B12-31-TLN
% Mức: TH
Khí nitrogen có $M=0,028\,\text{kg}$ /mol ở 300 K. Lấy $R$ = $8,31\,\text{J}$ /(mol·K). Tính tốc độ căn quân phương gần đúng.
\shortans{517}
\loigiai{
$v_{rms}=\sqrt{3RT/M}=\sqrt{3\cdot8,31\cdot300/0,028}\approx517$ m/s.
}
\end{ex}

% ===== Câu 107 =====
% ID: L12C2B12-31-TN
% Mức: NB
\dangbt{Nguồn gốc áp suất khí và va chạm phân tử}
\begin{ex}
% ID: L12C2B12-31-TN
% Mức: NB
Trong dạng “Nguồn gốc áp suất khí và va chạm phân tử”, Áp suất chất khí lên thành bình có nguồn gốc từ
\choice
{\True các va chạm của phân tử khí với thành bình}
{trọng lượng riêng của bình}
{lực hút của Trái Đất בלבד}
{thể tích của phân tử}
\loigiai{
Đáp án A. Phân tử đổi động lượng khi va chạm và gây lực lên thành bình.
}
\end{ex}

% ===== Câu 108 =====
% ID: L12C2B12-32-TL
% Mức: TH
\dangbt{Tính tốc độ căn quân phương của phân tử khí}
\begin{ex}
% ID: L12C2B12-32-TL
% Mức: TH
Nêu quy trình giải một bài thuộc dạng “Tính tốc độ căn quân phương của phân tử khí”.
\choiceTF

\loigiai{
Bước 1: tóm tắt và đổi đơn vị. Bước 2: nhận diện khí lí tưởng ở nhiệt độ tuyệt đối T. Bước 3: dùng $v_{rms}=\sqrt{3RT/M}=\sqrt{3kT/m}$. Bước 4: giải, ghi đơn vị và đối chiếu $v_{rms}$ tỉ lệ với $\sqrt{T/M}$.
}
\end{ex}

% ===== Câu 109 =====
% ID: L12C2B12-32-TLN
% Mức: TH
\begin{ex}
% ID: L12C2B12-32-TLN
% Mức: TH
Đổi 27 ° $C$ sang nhiệt độ tuyệt đối theo kelvin.
\shortans{300}
\loigiai{
$T=27+273=300$ K.
}
\end{ex}

% ===== Câu 110 =====
% ID: L12C2B12-32-TN
% Mức: NB
\dangbt{Nguồn gốc áp suất khí và va chạm phân tử}
\begin{ex}
% ID: L12C2B12-32-TN
% Mức: NB
Trong dạng “Nguồn gốc áp suất khí và va chạm phân tử”, Công thức áp suất theo mô hình động học phân tử là
\choice
{\True $p=\frac13\rho\overline{v^2}$}
{$p=\rho gh$}
{$p=\rho/3v^2$}
{$p=3\rho\overline{v^2}$}
\loigiai{
Đáp án A. Mô hình động học cho $p=\frac13\rho\overline{v^2}$.
}
\end{ex}

% ===== Câu 111 =====
% ID: L12C2B12-33-TL
% Mức: VD
\dangbt{Tính tốc độ căn quân phương của phân tử khí}
\begin{ex}
% ID: L12C2B12-33-TL
% Mức: VD
Khí nitrogen có $M=0,028\,\text{kg}$ /mol ở 300 K. Lấy $R$ = $8,31\,\text{J}$ /(mol·K). Tính tốc độ căn quân phương gần đúng.
\choiceTF

\loigiai{
$v_{rms}=\sqrt{3RT/M}=\sqrt{3\cdot8,31\cdot300/0,028}\approx517$ m/s.
}
\end{ex}

% ===== Câu 112 =====
% ID: L12C2B12-33-TLN
% Mức: VD
\begin{ex}
% ID: L12C2B12-33-TLN
% Mức: VD
Một đại lượng tăng từ 100 lên 150. Tính tỉ số giá trị sau so với trước.
\shortans{1,5}
\loigiai{
Tỉ số bằng $150/100=1,5$.
}
\end{ex}

% ===== Câu 113 =====
% ID: L12C2B12-33-TN
% Mức: TH
\dangbt{Nguồn gốc áp suất khí và va chạm phân tử}
\begin{ex}
% ID: L12C2B12-33-TN
% Mức: TH
Trong dạng “Nguồn gốc áp suất khí và va chạm phân tử”, Khí có $\rho=1,3$ kg/m³ và tốc độ căn quân phương $460\,\text{m}$ /s. Áp suất gần bằng
\choice
{\True $91,69\,\text{kPa}$}
{$275,08\,\text{kPa}$}
{$30,56\,\text{kPa}$}
{$0,6\,\text{kPa}$}
\loigiai{
Đáp án A. $p=\rho v_{rms}^2/3=91,69$ kPa.
}
\end{ex}

% ===== Câu 114 =====
% ID: L12C2B12-34-TL
% Mức: VD
\dangbt{Tính tốc độ căn quân phương của phân tử khí}
\begin{ex}
% ID: L12C2B12-34-TL
% Mức: VD
Giải thích vì sao cần tránh sai lầm: “dùng khối lượng mol theo g/mol mà chưa đổi sang kg/mol”.
\choiceTF

\loigiai{
Sai lầm này làm dữ kiện không cùng hệ đơn vị hoặc không đúng bản chất đại lượng, khiến phép thế vào $v_{rms}=\sqrt{3RT/M}=\sqrt{3kT/m}$ cho kết quả sai. Cần chuẩn hóa trước khi tính.
}
\end{ex}

% ===== Câu 115 =====
% ID: L12C2B12-34-TLN
% Mức: VD
\begin{ex}
% ID: L12C2B12-34-TLN
% Mức: VD
Một thể tích 2,5 $L$ bằng bao nhiêu m³? Viết kết quả theo dạng $a\cdot10^{-3}$.
\shortans{2,5}
\loigiai{
$2,5L=2,5\cdot10^{-3}$ m³.
}
\end{ex}

% ===== Câu 116 =====
% ID: L12C2B12-34-TN
% Mức: TH
\dangbt{Nguồn gốc áp suất khí và va chạm phân tử}
\begin{ex}
% ID: L12C2B12-34-TN
% Mức: TH
Trong dạng “Nguồn gốc áp suất khí và va chạm phân tử”, Khi nhiệt độ tuyệt đối tăng gấp 4, tốc độ căn quân phương của phân tử
\choice
{\True tăng gấp 2}
{tăng gấp 4}
{giảm một nửa}
{không đổi}
\loigiai{
Đáp án A. $v_{rms}\sim\sqrt{T}$.
}
\end{ex}

% ===== Câu 117 =====
% ID: L12C2B12-35-TL
% Mức: VD
\dangbt{Tính tốc độ căn quân phương của phân tử khí}
\begin{ex}
% ID: L12C2B12-35-TL
% Mức: VD
Phân tích cách kiểm tra tính hợp lí của kết quả trong dạng “Tính tốc độ căn quân phương của phân tử khí”.
\choiceTF

\loigiai{
Kiểm tra thứ nguyên theo m/s; kiểm tra dấu và bậc lớn; cuối cùng đối chiếu xu hướng $v_{rms}$ tỉ lệ với $\sqrt{T/M}$. Nếu trái xu hướng phải xem lại điều kiện và phép đổi đơn vị.
}
\end{ex}

% ===== Câu 118 =====
% ID: L12C2B12-35-TLN
% Mức: VD
\begin{ex}
% ID: L12C2B12-35-TLN
% Mức: VD
Áp suất $120\,\text{kPa}$ bằng bao nhiêu Pa, lấy hệ số của $10^3$ làm đáp án?
\shortans{120}
\loigiai{
$120\,\text{kPa}=120\cdot10^3$ Pa.
}
\end{ex}

% ===== Câu 119 =====
% ID: L12C2B12-35-TN
% Mức: TH
\dangbt{Nguồn gốc áp suất khí và va chạm phân tử}
\begin{ex}
% ID: L12C2B12-35-TN
% Mức: TH
Trong dạng “Nguồn gốc áp suất khí và va chạm phân tử”, Động năng tịnh tiến trung bình của một phân tử khí lí tưởng bằng
\choice
{\True $\frac32kT$}
{$kT$}
{$\frac12kT$}
{$3kT$}
\loigiai{
Đáp án A. Theo thuyết động học, $\overline{E_k}=3kT/2$.
}
\end{ex}

% ===== Câu 120 =====
% ID: L12C2B12-36-TL
% Mức: VDC
\dangbt{Tính tốc độ căn quân phương của phân tử khí}
\begin{ex}
% ID: L12C2B12-36-TL
% Mức: VDC
Đề xuất cách giải khi bài toán có nhiều giai đoạn liên quan đến dạng “Tính tốc độ căn quân phương của phân tử khí”.
\choiceTF

\loigiai{
Chia quá trình thành các trạng thái liên tiếp, ghi p, $V$, $T$ cho từng trạng thái; áp dụng $v_{rms}=\sqrt{3RT/M}=\sqrt{3kT/m}$ ở đoạn thỏa khí lí tưởng ở nhiệt độ tuyệt đối $T$; dùng kết quả đoạn trước làm dữ kiện đoạn sau và kiểm tra tính liên tục.
}
\end{ex}

% ===== Câu 121 =====
% ID: L12C2B12-36-TLN
% Mức: VDC
\begin{ex}
% ID: L12C2B12-36-TLN
% Mức: VDC
Theo quan hệ “ $v_{rms}$ tỉ lệ với $\sqrt{T/M}$ ”, nếu đại lượng phụ thuộc tỉ lệ thuận tăng gấp 4 thì hệ số tăng là bao nhiêu?
\shortans{4}
\loigiai{
Với quan hệ tỉ lệ thuận, hệ số biến đổi được giữ nguyên: 4.
}
\end{ex}

% ===== Câu 122 =====
% ID: L12C2B12-36-TN
% Mức: VD
\dangbt{Nguồn gốc áp suất khí và va chạm phân tử}
\begin{ex}
% ID: L12C2B12-36-TN
% Mức: VD
Trong dạng “Nguồn gốc áp suất khí và va chạm phân tử”, Hệ thức liên hệ áp suất với mật độ phân tử n là
\choice
{\True $p=nkT$}
{$p=nRT$}
{$p=k/(nT)$}
{$p=nk/T$}
\loigiai{
Đáp án A. Với n là số phân tử trên một đơn vị thể tích, $p=nkT$.
}
\end{ex}

% ===== Câu 123 =====
% ID: L12C2B12-37-TN
% Mức: VD
\begin{ex}
% ID: L12C2B12-37-TN
% Mức: VD
Trong dạng “Nguồn gốc áp suất khí và va chạm phân tử”, Ở cùng nhiệt độ, phân tử khí có khối lượng nhỏ hơn có tốc độ căn quân phương
\choice
{\True lớn hơn}
{nhỏ hơn}
{bằng nhau}
{bằng không}
\loigiai{
Đáp án A. $v_{rms}=\sqrt{3kT/m}$.
}
\end{ex}

% ===== Câu 124 =====
% ID: L12C2B12-38-TN
% Mức: VD
\begin{ex}
% ID: L12C2B12-38-TN
% Mức: VD
Trong dạng “Nguồn gốc áp suất khí và va chạm phân tử”, Nếu giữ mật độ phân tử không đổi và tăng $T$ gấp đôi thì áp suất
\choice
{\True tăng gấp đôi}
{giảm một nửa}
{không đổi}
{tăng gấp bốn}
\loigiai{
Đáp án A. Từ $p=nkT$, p tỉ lệ thuận T.
}
\end{ex}

% ===== Câu 125 =====
% ID: L12C2B12-39-TN
% Mức: VDC
\begin{ex}
% ID: L12C2B12-39-TN
% Mức: VDC
Trong dạng “Nguồn gốc áp suất khí và va chạm phân tử”, Khi va chạm đàn hồi với thành bình đứng yên, đại lượng của phân tử đổi dấu là
\choice
{\True thành phần vận tốc vuông góc thành}
{khối lượng}
{động năng}
{độ lớn vận tốc}
\loigiai{
Đáp án A. Thành phần pháp tuyến đảo chiều; tốc độ và động năng được bảo toàn.
}
\end{ex}

% ===== Câu 126 =====
% ID: L12C2B12-40-TN
% Mức: VDC
\begin{ex}
% ID: L12C2B12-40-TN
% Mức: VDC
Trong dạng “Nguồn gốc áp suất khí và va chạm phân tử”, Đơn vị của hằng số Boltzmann k là
\choice
{\True J/ $K$}
{J· $K$}
{Pa/ $K$}
{mol/ $J$}
\loigiai{
Đáp án A. Từ $E=kT$, k có đơn vị J/K.
}
\end{ex}

% ===== Câu 127 =====
% ID: L12C2B12-41-TN
% Mức: NB
\dangbt{So sánh chuyển động nhiệt của các loại khí}
\begin{ex}
% ID: L12C2B12-41-TN
% Mức: NB
Công thức cốt lõi phù hợp nhất cho dạng “So sánh chuyển động nhiệt của các loại khí” là
\choice
{\True $v_{rms}=\sqrt{3RT/M}$}
{$Q=mc\Delta T$}
{$A=Fs$}
{$U=RI$}
\loigiai{
Đáp án A. Dùng $v_{rms}=\sqrt{3RT/M}$ trong điều kiện các khí lí tưởng được so sánh ở nhiệt độ xác định.
}
\end{ex}

% ===== Câu 128 =====
% ID: L12C2B12-42-TN
% Mức: NB
\begin{ex}
% ID: L12C2B12-42-TN
% Mức: NB
Điều kiện quan trọng khi áp dụng dạng “So sánh chuyển động nhiệt của các loại khí” là
\choice
{khối lượng vật luôn bằng $1\,\text{kg}$}
{\True các khí lí tưởng được so sánh ở nhiệt độ xác định}
{nhiệt độ luôn bằng 0 ° $C$}
{áp suất luôn bằng áp suất khí quyển}
\loigiai{
Đáp án B. Điều kiện áp dụng là các khí lí tưởng được so sánh ở nhiệt độ xác định.
}
\end{ex}

% ===== Câu 129 =====
% ID: L12C2B12-43-TN
% Mức: TH
\begin{ex}
% ID: L12C2B12-43-TN
% Mức: TH
Nhận xét đúng về quan hệ các đại lượng là
\choice
{các đại lượng không phụ thuộc nhau}
{mọi đại lượng đều không đổi}
{\True ở cùng $T$, $v_{rms}\sim1/\sqrt M$}
{chỉ phụ thuộc thời gian}
\loigiai{
Đáp án C. ở cùng $T$, $v_{rms}\sim1/\sqrt M$.
}
\end{ex}

% ===== Câu 130 =====
% ID: L12C2B12-44-TN
% Mức: TH
\begin{ex}
% ID: L12C2B12-44-TN
% Mức: TH
Sai lầm cần tránh khi giải dạng này là
\choice
{ghi đầy đủ dữ kiện}
{đổi đơn vị đúng}
{kiểm tra điều kiện}
{\True cho rằng mọi loại phân tử ở cùng nhiệt độ có cùng tốc độ}
\loigiai{
Đáp án D. Cần tránh: cho rằng mọi loại phân tử ở cùng nhiệt độ có cùng tốc độ.
}
\end{ex}

% ===== Câu 131 =====
% ID: L12C2B12-45-TN
% Mức: TH
\begin{ex}
% ID: L12C2B12-45-TN
% Mức: TH
Đơn vị thường dùng sau khi chuẩn hóa theo SI là
\choice
{\True m/s}
{kg·m}
{W/s}
{C/mol}
\loigiai{
Đáp án A. Các đơn vị phù hợp: m/s.
}
\end{ex}

% ===== Câu 132 =====
% ID: L12C2B12-46-TN
% Mức: VD
\begin{ex}
% ID: L12C2B12-46-TN
% Mức: VD
Bước đầu tiên hợp lí khi giải bài là
\choice
{thay số ngay}
{\True xác định quá trình và đại lượng không đổi}
{bỏ qua đơn vị}
{đổi áp suất thành nhiệt độ}
\loigiai{
Đáp án B. Trước hết phải nhận diện quá trình, liệt kê dữ kiện và đại lượng không đổi.
}
\end{ex}

% ===== Câu 133 =====
% ID: L12C2B12-47-TN
% Mức: VD
\begin{ex}
% ID: L12C2B12-47-TN
% Mức: VD
Với dữ kiện của bài mẫu: Ở cùng nhiệt độ, khối lượng mol của khí $B$ gấp 4 lần khí A. Tính tỉ số $v_A/v_B$.
\choice
{\True Kết quả 2}
{Kết quả 4}
{Không đủ dữ kiện}
{Kết quả bằng 0}
\loigiai{
Đáp án A. Ở cùng $T$, $v\sim1/\sqrt M$, nên $v_A/v_B=\sqrt{M_B/M_A}=2$.
}
\end{ex}

% ===== Câu 134 =====
% ID: L12C2B12-48-TN
% Mức: VD
\begin{ex}
% ID: L12C2B12-48-TN
% Mức: VD
Khi kiểm tra kết quả, tiêu chí vật lí quan trọng là
\choice
{chỉ cần số dương}
{không cần đơn vị}
{\True phù hợp chiều biến thiên theo định luật}
{luôn phải lớn hơn giá trị đầu}
\loigiai{
Đáp án C. Kết quả phải phù hợp với nhận xét: ở cùng $T$, $v_{rms}\sim1/\sqrt M$.
}
\end{ex}

% ===== Câu 135 =====
% ID: L12C2B12-49-TN
% Mức: VDC
\begin{ex}
% ID: L12C2B12-49-TN
% Mức: VDC
Nếu giữ đúng các điều kiện của định luật, biểu thức nên dùng là
\choice
{$s=vt$}
{\True $v_{rms}=\sqrt{3RT/M}$}
{$P=A/t$}
{$F=ma$}
\loigiai{
Đáp án B. Biểu thức đặc trưng là $v_{rms}=\sqrt{3RT/M}$.
}
\end{ex}

% ===== Câu 136 =====
% ID: L12C2B12-50-TN
% Mức: VDC
\begin{ex}
% ID: L12C2B12-50-TN
% Mức: VDC
Phát biểu nào phản ánh đúng bản chất bài toán?
\choice
{Có thể bỏ qua nhiệt độ tuyệt đối}
{Không cần xác định lượng khí}
{Mọi quá trình đều đẳng nhiệt}
{\True Phải dùng $v_{rms}=\sqrt{3RT/M}$ và kiểm tra các khí lí tưởng được so sánh ở nhiệt độ xác định}
\loigiai{
Đáp án D. Phải dùng $v_{rms}=\sqrt{3RT/M}$, đồng thời kiểm tra các khí lí tưởng được so sánh ở nhiệt độ xác định.
}
\end{ex}

% ===== Câu 137 =====
% ID: L12C2B12-51-TN
% Mức: NB
\dangbt{Tính tốc độ căn quân phương của phân tử khí}
\begin{ex}
% ID: L12C2B12-51-TN
% Mức: NB
Công thức cốt lõi phù hợp nhất cho dạng “Tính tốc độ căn quân phương của phân tử khí” là
\choice
{\True $v_{rms}=\sqrt{3RT/M}=\sqrt{3kT/m}$}
{$Q=mc\Delta T$}
{$A=Fs$}
{$U=RI$}
\loigiai{
Đáp án A. Dùng $v_{rms}=\sqrt{3RT/M}=\sqrt{3kT/m}$ trong điều kiện khí lí tưởng ở nhiệt độ tuyệt đối T.
}
\end{ex}

% ===== Câu 138 =====
% ID: L12C2B12-52-TN
% Mức: NB
\begin{ex}
% ID: L12C2B12-52-TN
% Mức: NB
Điều kiện quan trọng khi áp dụng dạng “Tính tốc độ căn quân phương của phân tử khí” là
\choice
{khối lượng vật luôn bằng $1\,\text{kg}$}
{\True khí lí tưởng ở nhiệt độ tuyệt đối $T$}
{nhiệt độ luôn bằng 0 ° $C$}
{áp suất luôn bằng áp suất khí quyển}
\loigiai{
Đáp án B. Điều kiện áp dụng là khí lí tưởng ở nhiệt độ tuyệt đối T.
}
\end{ex}

% ===== Câu 139 =====
% ID: L12C2B12-53-TN
% Mức: TH
\begin{ex}
% ID: L12C2B12-53-TN
% Mức: TH
Nhận xét đúng về quan hệ các đại lượng là
\choice
{các đại lượng không phụ thuộc nhau}
{mọi đại lượng đều không đổi}
{\True $v_{rms}$ tỉ lệ với $\sqrt{T/M}$}
{chỉ phụ thuộc thời gian}
\loigiai{
Đáp án C. $v_{rms}$ tỉ lệ với $\sqrt{T/M}$.
}
\end{ex}

% ===== Câu 140 =====
% ID: L12C2B12-54-TN
% Mức: TH
\begin{ex}
% ID: L12C2B12-54-TN
% Mức: TH
Sai lầm cần tránh khi giải dạng này là
\choice
{ghi đầy đủ dữ kiện}
{đổi đơn vị đúng}
{kiểm tra điều kiện}
{\True dùng khối lượng mol theo g/mol mà chưa đổi sang kg/mol}
\loigiai{
Đáp án D. Cần tránh: dùng khối lượng mol theo g/mol mà chưa đổi sang kg/mol.
}
\end{ex}

% ===== Câu 141 =====
% ID: L12C2B12-55-TN
% Mức: TH
\begin{ex}
% ID: L12C2B12-55-TN
% Mức: TH
Đơn vị thường dùng sau khi chuẩn hóa theo SI là
\choice
{\True m/s}
{kg·m}
{W/s}
{C/mol}
\loigiai{
Đáp án A. Các đơn vị phù hợp: m/s.
}
\end{ex}

% ===== Câu 142 =====
% ID: L12C2B12-56-TN
% Mức: VD
\begin{ex}
% ID: L12C2B12-56-TN
% Mức: VD
Bước đầu tiên hợp lí khi giải bài là
\choice
{thay số ngay}
{\True xác định quá trình và đại lượng không đổi}
{bỏ qua đơn vị}
{đổi áp suất thành nhiệt độ}
\loigiai{
Đáp án B. Trước hết phải nhận diện quá trình, liệt kê dữ kiện và đại lượng không đổi.
}
\end{ex}

% ===== Câu 143 =====
% ID: L12C2B12-57-TN
% Mức: VD
\begin{ex}
% ID: L12C2B12-57-TN
% Mức: VD
Với dữ kiện của bài mẫu: Khí nitrogen có $M=0,028\,\text{kg}$ /mol ở 300 K. Lấy $R$ = $8,31\,\text{J}$ /(mol·K). Tính tốc độ căn quân phương gần đúng.
\choice
{\True Kết quả 517}
{Kết quả 1034}
{Không đủ dữ kiện}
{Kết quả bằng 0}
\loigiai{
Đáp án A. $v_{rms}=\sqrt{3RT/M}=\sqrt{3\cdot8,31\cdot300/0,028}\approx517$ m/s.
}
\end{ex}

% ===== Câu 144 =====
% ID: L12C2B12-58-TN
% Mức: VD
\begin{ex}
% ID: L12C2B12-58-TN
% Mức: VD
Khi kiểm tra kết quả, tiêu chí vật lí quan trọng là
\choice
{chỉ cần số dương}
{không cần đơn vị}
{\True phù hợp chiều biến thiên theo định luật}
{luôn phải lớn hơn giá trị đầu}
\loigiai{
Đáp án C. Kết quả phải phù hợp với nhận xét: $v_{rms}$ tỉ lệ với $\sqrt{T/M}$.
}
\end{ex}

% ===== Câu 145 =====
% ID: L12C2B12-59-TN
% Mức: VDC
\begin{ex}
% ID: L12C2B12-59-TN
% Mức: VDC
Nếu giữ đúng các điều kiện của định luật, biểu thức nên dùng là
\choice
{$s=vt$}
{\True $v_{rms}=\sqrt{3RT/M}=\sqrt{3kT/m}$}
{$P=A/t$}
{$F=ma$}
\loigiai{
Đáp án B. Biểu thức đặc trưng là $v_{rms}=\sqrt{3RT/M}=\sqrt{3kT/m}$.
}
\end{ex}

% ===== Câu 146 =====
% ID: L12C2B12-60-TN
% Mức: VDC
\begin{ex}
% ID: L12C2B12-60-TN
% Mức: VDC
Phát biểu nào phản ánh đúng bản chất bài toán?
\choice
{Có thể bỏ qua nhiệt độ tuyệt đối}
{Không cần xác định lượng khí}
{Mọi quá trình đều đẳng nhiệt}
{\True Phải dùng $v_{rms}=\sqrt{3RT/M}=\sqrt{3kT/m}$ và kiểm tra khí lí tưởng ở nhiệt độ tuyệt đối $T$}
\loigiai{
Đáp án D. Phải dùng $v_{rms}=\sqrt{3RT/M}=\sqrt{3kT/m}$, đồng thời kiểm tra khí lí tưởng ở nhiệt độ tuyệt đối T.
}
\end{ex}
